\PassOptionsToPackage{dvipsnames}{xcolor}
\documentclass[acmsmall,screen,review,anonymous,nonacm]{acmart}
\citestyle{acmauthoryear}
\bibliographystyle{ACM-Reference-Format}

% Each paper should have no more than 25 pages of text, excluding bibliography

\usepackage{mathpartir}
\usepackage{stmaryrd}
\usepackage{placeins}
\usepackage{mathtools}
\usepackage{subcaption}
\usepackage{cleveref}
\usepackage{listings}
\usepackage{tikz}
\usepackage{xspace}
\usepackage{xr-hyper}
\usetikzlibrary{arrows.meta,calc,fit,positioning,shapes.multipart}

\crefname{lstlisting}{listing}{listings}
\Crefname{lstlisting}{Listing}{Listings}

\definecolor{verylightgray}{rgb}{0.9,0.9,0.9}
\definecolor{verylightred}{rgb}{1,0.9,0.9}
\definecolor{darkgold}{rgb}{0.7,0.5.4,0}
\definecolor{darkviolet}{rgb}{0.5,0,0.4}
\definecolor{darkgreen}{rgb}{0,0.4,0.2}
\definecolor{darkblue}{rgb}{0.1,0.1,0.9}
\definecolor{darkgrey}{rgb}{0.5,0.5,0.5}
\definecolor{lightblue}{rgb}{0.4,0.4,1}

\lstdefinestyle{eclipse}{
  breaklines=true,
  basicstyle=\sffamily\small,
  emphstyle=\color{red}\bfseries,
  keywordstyle=\color{darkviolet}\bfseries,
  commentstyle=\color{darkgreen},
  stringstyle=\color{darkblue},
  numberstyle=\color{darkgrey},%\lstfontfamily,
  emphstyle=\color{red},
  showstringspaces=false,
}

\lstdefinestyle{eclipse-new}{
  backgroundcolor=\color{verylightgray},
  breaklines=true,
  basicstyle=\sffamily\small,
  emphstyle=\color{red}\bfseries,
  keywordstyle=\color{darkgold}\bfseries,
  commentstyle=\color{darkgreen},
  stringstyle=\color{darkblue},
  numberstyle=\color{darkgrey},%\lstfontfamily,
  emphstyle=\color{red},
  showstringspaces=false,
}

\lstdefinestyle{eclipse-error}{
  backgroundcolor=\color{verylightred},
  breaklines=true,
  basicstyle=\sffamily\small,
  emphstyle=\color{red}\bfseries,
  keywordstyle=\color{darkgold}\bfseries,
  commentstyle=\color{darkgreen},
  stringstyle=\color{darkblue},
  numberstyle=\color{darkgrey},%\lstfontfamily,
  emphstyle=\color{red},
  showstringspaces=false,
}

\lstset{
  language=[Objective]Caml,
  style=eclipse,
  morekeywords=[1]{Int, Char, Bool, Str, String, Close, Wait, Return, Unit, Skip, Chan, Ret, void, type, dualof, rec, let, in, if, then,
    else, new, send, recv, select, fork, close, wait, print, branch, offer,
    pure, full, drop, case, of, data, match, with,
    read, write, free, split, send_borrow, recv_borrow, True,
    False}, %  MU, ML, SU, SL, TU, TL
  % sensitive=true,
  literate=
  {>>}{>>}1
  {|>}{$\rhd$}1
  {->}{$\rightarrow$}2
  {-R>}{${\rightarrow_{\textsc r}}$}3
  {-L>}{${\rightarrow_{\textsc l}}$}3
  {-o}{$\multimap$}2
  {.o}{$\odot$}2
  {ox}{$\otimes$}2
  {=>}{$\Rightarrow$}2
  {forall}{$\forall$}1
  {Lambda}{$\Lambda$}1
  {lambda}{$\lambda$}1
  {mu}{$\mu$}1
  {nu}{$\nu$}1
  {alpha}{$\alpha$}1
  {sigma}{$\sigma$}1
  {oplus}{$\oplus$}1
  {+\{}{$\oplus$\{}2
  {\\le}{$\le$}1
  {elll}{$\ell$}1
  {_1}{${}_1$}1
  {_n}{${}_n$}1
  ,
  morecomment=[l]{\#},
  breaklines=true,
  tabsize=2
}

%%% Local Variables:
%%% mode: latex
%%% TeX-master: "popl2027"
%%% End:

\newcommand{\thecalculus}{\textsc{BGV}\xspace}

\newcommand{\metacolor}{\textcolor{black}}
\newcommand{\ctxcolor}{\textcolor{RoyalPurple}}
\newcommand{\targetctxcolor}{\textcolor{blue}}
\newcommand{\typecolor}{\textcolor{RoyalPurple}}
\newcommand{\targettypecolor}{\textcolor{blue}}
\newcommand{\termcolor}{\textcolor{Mahogany}}
\newcommand{\targettermcolor}{\textcolor{Sepia}}
\newcommand{\proccolor}{\termcolor}
\newcommand{\targetproccolor}{\targettermcolor}
\newcommand{\procbindcolor}{\termcolor}
\newcommand{\targetprocbindcolor}{\targettermcolor}
\newcommand{\priocolor}{\textcolor{Emerald}}  % LimeGreen -> too light
\newcommand{\dircolor}{\textcolor{Melon}}
\newcommand{\mobcolor}{\textcolor{LimeGreen}}

\newcommand*\Keyword[1]{\operatorname{\texttt{#1}}}
\newcommand*\TermKeyword[1]{\termcolor{\Keyword{#1}}}
\newcommand*\TypeKeyword[1]{\typecolor{\Keyword{#1}}}
\newcommand*\TargetTypeKeyword[1]{\targettypecolor{\Keyword{#1}}}
\newcommand*\TargetTermKeyword[1]{\targettermcolor{\Keyword{#1}}}

\newcommand*\NewMetaVar[3][]{%
  \NewDocumentCommand #2 {s s O{}} {%
    #1{#3\IfBooleanT{##1}{\IfBooleanTF{##2}{''}{'}}_{##3}}%
  }%
}
\newcommand*\NewContextApp[3][]{%
  \NewDocumentCommand #2 {s O{} m} {%
    #1{\IfBooleanTF{##1}{#3*[##2]}{#3[##2]}[##3]}%
  }%
}

\DeclareMathOperator\dom{\mathsf{dom}}

% Grammars
\newcommand{\grmeq}{\; \Coloneqq \;\;}
\newcommand{\grmor}{\;\mid\;}

% Mobilities
\newcommand*\Mob{\mobcolor{m}}
\newcommand*\MobMobile{\mobcolor{\textsc m}}
\newcommand*\MobStationary{\mobcolor{\textsc s}}

% Multiplicities
\newcommand*\Mul{m}
\newcommand*\MulLin{1}
\newcommand*\MulUnr{\omega}

% Directions
\newcommand*\Dir{\dircolor{d}}
\newcommand*\DirLeft{\dircolor{\textsc l}}
\newcommand*\DirRight{\dircolor{\textsc r}}
\newcommand*\DirNone{\dircolor{\upomega}} % {\cdot} % This should be un
\newcommand*\DirUnord{\dircolor{\mathbf 1}} % {\circ} %This should be lin

% Pairs
\newcommand*\Pair[1][{}]{\mathbin{\typecolor{\otimes}_{#1}}}
\newcommand*\PairUnord{\mathbin{\typecolor{\otimes}}}
\newcommand*\PairOrd{\mathbin{\typecolor{\odot}}}

% Kinds
\newcommand*\K{K}
\newcommand*\KSession[1]{\mathcal{S}^{#1}}
\newcommand*\KType[1]{\mathcal{T}^{#1}}

% Base Types
\newcommand*\B{\typecolor{B}}
\newcommand*\BUnit{\TypeKeyword{Unit}}
\newcommand*\BInt{\TypeKeyword{Int}}
\newcommand*\BBool{\TypeKeyword{Bool}}
\newcommand*\BString{\TypeKeyword{String}}

% Types
\NewMetaVar[\typecolor] \T T
\newcommand*\TProd[3]{\typecolor{{#2}\otimes_{#1}{#3}}}
\newcommand*\TPair[3][{}]{\typecolor{{#2}\Pair[#1]{#3}}}
\newcommand*\TPairUnord[2]{\TProd\DirUnord{#1}{#2}} % \typecolor{{#1}\PairUnord{#2}}}
\newcommand*\TPairOrd[2]{\TProd\DirLeft{#1}{#2}} %{\typecolor{{#1}\PairOrd{#2}}}
\newcommand*\TSum[2]{\typecolor{{#1}+{#2}}}
\newcommand*\TVariant[2]{\typecolor{\langle{#1}\rangle_{#2}}}
\newcommand*\TArr[4]{\typecolor{{#3}\to^{#1}_{#2}{#4}}}
% \newcommand*\TArr[4]{{#3}\overset{{#1};{#2}}\to{#4}}
\newcommand*\TExists[3][{}]{\typecolor{\exists #2^{#1}. #3}}

% Unrestricted Types
\newcommand*\M{M}

% Session Types
\let\S=\undefined
\NewMetaVar[\typecolor] \S S
\NewMetaVar[\typecolor] \R R
\newcommand*\SVar{\typecolor{X}}
\newcommand*\SSkip{\TypeKeyword{Skip}}
\newcommand*\SSeq[2]{\typecolor{{#1};{#2}}}
\newcommand*\SWait[1][{}]{\TypeKeyword{Wait}}
\newcommand*\STerm[1][{}]{\TypeKeyword{Close}} % Term}}
\newcommand*\SDual[1]{\typecolor{#1}^{\metacolor\bot}} % {\typecolor{\Dual\,({#1})}}
\newcommand*\STrans[2][{}]{\typecolor{\operatorname{\#}{#2}}}
\newcommand*\SSend[1]{\typecolor{\operatorname{!}{\!#1}}}
\newcommand*\SRecv[1]{\typecolor{\operatorname{?}{\!#1}}}
\newcommand*\SRecord[2]{\typecolor{\{{#1}\}_{#2}}}
\newcommand*\SVariant[2]{\typecolor{\langle{#1}\rangle_{#2}}}
\newcommand*\SMulti[3][{}]{\typecolor{\odot\SRecord{#2}{#3}}}
\newcommand*\SChoice[3][{}]{\typecolor{\oplus\SRecord{#2}{#3}}}
\newcommand*\SBranch[3][{}]{\typecolor{\&\SRecord{#2}{#3}}}
\newcommand*\SRec[2]{\typecolor{\mu{#1}.{#2}}}
\newcommand*\SRet{\TypeKeyword{Ret}}
\newcommand*\SAcq{\TypeKeyword{Acq}}

\newcommand*\SComp[1]{{#1}\mathbin{\metacolor{\fatsemi}}}

\newcommand\SSeen{\ctxcolor{\Theta}}
\newcommand\SeenEmpty{\ctxcolor{\emptyset}}
\newcommand\SeenSnoc[1]{\ctxcolor{;{}}({#1})}

\newcommand*\Sbot[1][{}]{\typecolor{S^{?}_{#1}}}
\newcommand\SNil{\typecolor{\bot}}
\newcommand\Ssub[2]{[{#2}/{#1}]}

\newcommand\SStrip[3][\SSeen]{{#1}\vdash{#2}\mathbin{\%}{#3}}
\DeclareMathOperator\Smerge{merge}

% Target
\NewMetaVar[\targettypecolor] \TT T
\NewMetaVar[\targettypecolor] \ST S
\newcommand*\STSend[2][{}]{\targettypecolor{\operatorname{!}^{\priocolor{#1}}{#2}.}}
\newcommand*\STRecv[2][{}]{\targettypecolor{\operatorname{?}^{\priocolor{#1}}{#2}.}}
\newcommand*\STWait[1][{}]{\targettypecolor{\Keyword{Wait}}^{\priocolor{#1}}}
\newcommand*\STTerm[1][{}]{\targettypecolor{\Keyword{Close}}^{\priocolor{#1}}} % Term}}
\newcommand*\STDual[1]{\targettypecolor{\Dual\,({#1})}}
\newcommand*\TBUnit{\TargetTypeKeyword{Unit}}
\newcommand*\TBChan{\TargetTypeKeyword{Chan}}
\newcommand*\TBChanD{\TargetTypeKeyword{DChan}}
\NewDocumentCommand\TBChanF{s}{\targettypecolor{\TargetTypeKeyword{FChan}\IfBooleanT{#1}{{}^?}}}

\newcommand*\TTVariant[2]{\targettypecolor{\langle{#1}\rangle_{#2}}}
\newcommand*\TTArr[2]{\targettypecolor{{#1}\to{#2}}}
\newcommand*\TTProd[2]{\targettypecolor{{#1}\times{#2}}}
\newcommand*\TTSum[2]{\targettypecolor{{#1}+{#2}}}

% Contexts
\NewMetaVar[\ctxcolor] \Ctx \Gamma
\newcommand*\CtxEmpty{\ctxcolor{\emptyset}}
% \newcommand*\CtxEmpty{\ctxcolor{\cdot}}
\newcommand*\CtxVar[2]{\ctxcolor{{\termcolor{#1}}:{\typecolor{#2}}}}
\newcommand*\CtxExt[2]{\ctxcolor{{#1},{#2}}}
\newcommand*\CtxExtUnord[2]{\ctxcolor{{#1} \parallel {#2}}}
\newcommand*\CtxPat[2]{\ctxcolor{\Ctx[#1] \CtxPatOp \Ctx[#2]}}
\newcommand*\CtxPatOp{*}
\newcommand*\CtxPatternRaw{\ctxcolor{\mathcal{G}}}
\newcommand*\CtxPattern[2][{}]{\ctxcolor{\mathcal{G}_{#1}[#2]}}
\newcommand*\CtxLeftPattern[1]{\ctxcolor{\mathcal{L}[#1]}}
\newcommand*\CtxExtInactive[4]{{#1},{#2}:^{#4}[{#3}]}
\newcommand*\CtxSeq[2]{\ctxcolor{{#1},{#2}}}
\newcommand*\CtxPar[2]{\ctxcolor{{#1}\parallel{#2}}}
\newcommand*\CtxSplit[4]{{#1}=^{#4}{#2}\fatsemi{#3}}
\newcommand*\CtxExtTarget[2]{{#1},{#2}}
% \newcommand*\CtxExtTarget[2]{{#1}\rhd{#2}}
\newcommand*\IsLeftPat{\operatorname{\textsf{LeftPat}}}
\newcommand*\CtxRestrict[2]{{#1}\downarrow{#2}}
\newcommand*\SubCtx{\preccurlyeq}
\newcommand*\CtxBind[2]{\ctxcolor{{#1}:{#2}}}

% Session Contexts
\newcommand*\SCtx{\Xi}
\newcommand*\SCtxEmpty{\cdot}
\newcommand*\SCtxExt[3]{{#1},{#2}:{#3}}

% binders
\NewMetaVar[\procbindcolor] \BindGroup B
\NewMetaVar[\procbindcolor] \Bind      b
\newcommand*\BSeq[2]{\procbindcolor{#1; #2}}
\newcommand*\BPar[2]{\procbindcolor{#1\|#2}}
\newcommand*\BEmp{\procbindcolor{\cdot}}
\DeclareMathOperator\BindHead{head}

% Expressions
\NewMetaVar[\termcolor] \E     e
\NewMetaVar[\termcolor] \EVar  x
\NewMetaVar[\termcolor] \EVary y
\newcommand*\EBorrow[1]{\termcolor{\&{#1}}}
\newcommand*\EMu[2][f\!]{\termcolor{\mu {#1}.{#2}}}
\newcommand*\EAbs[3][{}]{\termcolor{\lambda_{#1}{#2}.{#3}}}
\newcommand*\EApp[2]{\termcolor{{#1}\,{#2}}}
\newcommand*\EAppD[3][{}]{\termcolor{({#2}\,{#3})_{#1}}} % with direction
\newcommand*\EInj[2]{\termcolor{\Keyword{inj$_{#1}$}{#2}}}
\newcommand*\EPair[3][{}]{\termcolor{{#2} \otimes_{#1} {#3}}}
\newcommand*\EPairOrd[2]{\EPair[\DirLeft]{#1}{#2}} % {\termcolor{{#1} \odot {#2}}}
\newcommand*\ECase[5]{\termcolor{\Keyword{case} {#1} \mathop{\{} \Keyword{inj$_1$} {#2} \Rightarrow{#3}; \Keyword{inj$_2$} {#4} \Rightarrow{#5} \mathop{\}}}}
\newcommand*\ELet[2]{\termcolor{\Keyword{let}\,{#1} \Keyword{=} {#2}\Keyword{in}\,}}
\newcommand*\ESeq[1]{\termcolor{{#1}\!\Keyword{;}}}
\newcommand*\ELetPair[4][{}]{\termcolor{\ELet{\EPair[#1]{#2}{#3}}{#4}}}
\newcommand*\ELetPairOrd[3]{\ELetPair[\DirLeft]{#1}{#2}{#3}} % {\termcolor{\ELet{\EPairOrd{#1}{#2}}{#3}}}
\newcommand*\EFork[1]{\termcolor{\CFork~{#1}}}
\newcommand*\EMatch[3]{\termcolor{\Keyword{match}\, {#1}\, \Keyword{with}\, \ERecord{#2}{#3}}}
\newcommand*\EMatchA[3]{\termcolor{\begin{array}[t]{@{}l@{}l}\Keyword{match}&\, {#1}\, \Keyword{with}\\&\ERecord{#2}{#3}\end{array}}}
\newcommand*\EPacked[2]{\termcolor{\langle#1,#2\rangle}}
\newcommand*\EPack[2]{\termcolor{\Keyword{pack}\,\langle#1,#2\rangle\,\Keyword{as}\,}}
\newcommand*\EOpen[3]{\termcolor{\Keyword{open}\,\langle#1,#2\rangle = #3\,\Keyword{in}\,}}
\newcommand*\ERecord[2]{\termcolor{\{{#1}\}_{#2}}}


\newcommand*\CRead{\targettermcolor{\mathit{read}}}
\newcommand*\CWrite{\targettermcolor{\mathit{write}}}
\newcommand*\CWrap{\targettermcolor{\mathit{wrap}}}

% Constants
\newcommand*\C[1][]{\termcolor{c#1}}
\newcommand*\CUnit{\TermKeyword{*}}
\newcommand*\CBase{\TermKeyword{b}}
\newcommand*\CFork{\TermKeyword{fork}}
\newcommand*\CSplit{\TermKeyword{lsplit}}
\newcommand*\CSSplit{\TermKeyword{rsplit}}
\newcommand*\CDrop{\TermKeyword{drop}}
\newcommand*\CAcq{\TermKeyword{acquire}}
\newcommand*\CNew{\TermKeyword{new}}
\newcommand*\CLink{\TermKeyword{link}}
\newcommand*\CWait{\TermKeyword{wait}}
\newcommand*\CTerm{\TermKeyword{close}} % term}}
\newcommand*\CSend{\TermKeyword{send}}
\newcommand*\CRecv{\TermKeyword{recv}} % recv}}
\newcommand*\CSelect[1]{\EApp{\TermKeyword{select}}{#1}}
\newcommand*\CBranch{\TermKeyword{branch}}
\newcommand*\COffer{\TermKeyword{offer}}


% target expressions
\NewMetaVar[\targettermcolor] \TEVar  x
\NewMetaVar[\targettermcolor] \TEVary y
\NewMetaVar[\targettermcolor] \TEVarz z
\NewDocumentCommand \TEChan {O{\CUnit} m O{\CUnit}}
  {
    \EPair {\EPair {#1} {#2}} {#3}
  }
    
% target contexts
\NewMetaVar[\targetctxcolor] \TCtx \Gamma
\newcommand*\TCtxEmpty{\targetctxcolor{\emptyset}}
\newcommand*\TCtxVar[2]{\targetctxcolor{{\targettermcolor{#1}}:{\targettypecolor{#2}}}}
\newcommand*\TCtxExt[2]{\targetctxcolor{{#1},{#2}}}
\newcommand*\TCtxApp[2]{\metacolor{\targetctxcolor{#1}\cup\targetctxcolor{#2}}}

\NewMetaVar[\targetprocbindcolor] \TBind B
\newcommand*\TBEmp{\targetprocbindcolor{\cdot}}
\newcommand*\TBExt[2]{\targetprocbindcolor{{#1},{#2}}}

% Substitutions
\NewMetaVar \Sub \sigma
\newcommand*\SubId{\mathrm{id}}
\newcommand*\SubExt[3][\Sub]{{#1}[{#2} \mapsto {#3}]}
\newcommand*\SubSingle[2]{\{{#1}\mapsto{#2}\}}

% Values
\NewMetaVar[\termcolor] \V v
\NewMetaVar[\termcolor] \U u
\NewMetaVar[\targettermcolor] \TV v
\NewMetaVar[\targettermcolor] \TU u

% Evaluation Contexts
\NewMetaVar[\termcolor] \EC E
\NewContextApp[\termcolor] \ECapp \EC
\newcommand*\ECEmpty{\termcolor{\Box}}

% Flagged Contexts
\NewMetaVar[\proccolor] \FC {\mathcal{F}}
\NewContextApp[\proccolor] \FCapp \FC

% Target Flagged Contexts
\NewMetaVar[\targetproccolor] \TFC {\mathcal{F}}
\NewContextApp[\targetproccolor] \TFCapp \TFC

% Configuration Contexts
\NewMetaVar[\proccolor] \CC {\mathcal{G}}
\newcommand*\CCapp[2][]{\proccolor{\CC[#1][#2]}}

% Target Configuration Contexts
\NewMetaVar[\targetproccolor] \TCC {\mathcal{G}}
\newcommand*\TCCapp[2][]{\targetproccolor{\TCC[#1][#2]}}
\newcommand*\TCCEmpty{\targetproccolor{\Box}}

% Processes
\NewMetaVar[\proccolor] \Proc P
\newcommand*\PExp[2][{}]{\proccolor{{}^{#1}\langle{#2}\rangle}}
\newcommand*\PNew[3]{\proccolor{(\nu [{#1}][{#2}]){#3}}}
\newcommand*\PPar[2]{\proccolor{{#1} \,\|\, {#2}}}

% Target Processes
\NewMetaVar[\targetproccolor] \TProc P
\newcommand*\TPExp[1]{\targetproccolor{\langle{#1}\rangle}}
\newcommand*\TPNew[2][xy]{\targetproccolor{(\TBindNu #1){#2}}}
\newcommand*\TPPar[2]{\targetproccolor{{#1} \,\|\, {#2}}}
\NewDocumentCommand\TPFlag{m m}{\targetproccolor{#1 \mapsto #2}}
\NewDocumentCommand\TPSync{s m m}{\targetproccolor{(\TBindSync[#1]#2){#3}}}

\newcommand*\TBindNu{\nu}
\NewDocumentCommand\TBindSync{s O{#1}}{\IfBooleanTF{#2}{\varphi^*}{\varphi}}

\NewMetaVar[\targetproccolor] \Flag \phi
\newcommand*\FlagSet{\targetproccolor 1}
\newcommand*\FlagUnset{\targetproccolor 0}

\newcommand*\BindFlag[1]{\metacolor{\phi\ifstrempty{#1}{}{(#1)}}}

% Variable Sets ("Dependencies")
\newcommand*\D{\Delta}
\newcommand*\DVar{\delta}
\newcommand*\DSingle[1]{\{{#1}\}}
\newcommand*\DEmpty{\emptyset}
\newcommand*\DJoin{\cup}

% Effects
\NewMetaVar[\typecolor] \Eff \epsilon
\newcommand*\EffPure{\typecolor{\textsc{p}}}
\newcommand*\EffImpure{\typecolor{\textsc{i}}}
% \newcommand*\EffPure{\typecolor{\mathbf{0}}}
% \newcommand*\EffImpure{\typecolor{\mathbf{1}}}

% Effect Order
\newcommand*\EffLeq{\sqsubseteq}
\newcommand*\EffJoin{\metacolor{{}\sqcup{}}} 

% Graphs
\newcommand*\Graph{G}
\newcommand*\Node{n}

% Session Type Equivalence
\newcommand*\SEq{\approx}

% Free Variables
\newcommand*\Fv{\operatorname{\mathsf{fv}}}

% Domain of a type context
\newcommand*\CtxDom{\operatorname{\mathsf{dom}}}

% Dual Session Type
\newcommand*\Dual{{\operatorname{\mathsf{dual}}}}

% Ordered Types
\newcommand*\Ord{\operatorname{\mathsf{ord}}}

% Unrestricted Types
\newcommand*\Unr{\operatorname{\mathsf{unr}}}

% Terminated Session Types
\newcommand*\Bounded{\operatorname{\mathsf{bounded}}}
\newcommand*\Mobile{\operatorname{\mathsf{mobile}}}
\newcommand*\Skips{\operatorname{\mathsf{skips}}}

% Session Type Formation
\newcommand*\HasSKind[4]{{#1}\mid{#2}\vdash{#3}:{#4}}

% Regular Type Formation
\newcommand*\HasTKind[3]{{#1}\vdash{#2}:{#3}}

% Expression Typing
\NewDocumentCommand \HasType {O{\Eff} m m m}
  {
    \ctxcolor{#2} \vdash {#3} : {#4} \ifstrempty{#1}{}{\mid{#1}}
  }

% Constant Typing
\newcommand*\HasCType[2]{{#1}:{#2}}

% Process Typing
\newcommand*\ProcHasType[2]{\ctxcolor{#1} \vdash {\termcolor{#2}} }
\newcommand*\BindGivesCtx[3]{#1 \vdash #2 \leadsto #3}

% Type Context Formation
\newcommand*\CtxWf[1]{\vdash{#1}}

% Multiplicity Order
\newcommand*\MulLeq{\leq}

% Subtyping for Sessions
\newcommand*\IsSSubtypeOf{\operatorname{<:^s}}

% Subtyping for Types
\newcommand*\IsTSubtypeOf{\operatorname{<:^t}}

% Algorithmic Expression Typing
\newcommand*\TypeInfer[5]{{#1}\vdash{#2}\uparrow{#3}\mid{#4}\mid{#5}}
\newcommand*\TypeCheck[5]{{#1}\vdash{#2}\downarrow{#3}\mid{#4}\mid{#5}}

\newcommand*\TypeInferX[5][]{{#2}\vdash{#3}\uparrow^{#1}{#4}\dashv{#5}}
\newcommand*\TypeCheckX[5][]{{#2}\vdash{#3}\downarrow^{#1}{#4}\dashv{#5}}

\newcommand*\CtxDiff[2]{{#1} \operatorname{\backslash} {#2}}
\newcommand*\IfThenElse[3]{\Keyword{if}{#1}\Keyword{then}{#2}\Keyword{else}{#3}}

% Reduction 
\newcommand{\Reduces}{\rightarrow}
\newcommand{\IsRed}[2]{{#1}\;{\reduces}\;{#2}}

% Operators 
\newcommand{\Subs}[3]{\termcolor{{#1}[{#2}/{#3}]}}

% Congruence
\newcommand{\Congruent}{\equiv}

% Translations
\newcommand*\TrCmd[2]{\metacolor{\mathcal{T}_{#1}(#2)}}
\newcommand*\TrRel[1]{\leadsto_{#1}}
\newcommand\TrT{\TrCmd T}
\newcommand\TrCtx{\TrCmd \Gamma}
\newcommand\TrBG{\TrCmd B}

\newcommand*\TransProc[1][\Sub]{\stackrel{#1}{\TrRel{}}}
\newcommand*\TransBind[3][\E]{\stackrel{\targettermcolor{#1}}{\TrRel{}} #2 \mid #3}

\newcommand*\TA[2]{\metacolor{\llbracket #1;#2 \rrbracket_{\Gamma}}}
% \newcommand*\TA[2]{\metacolor{\llbracket #2 \rrbracket_{\Gamma} ({#1})}}
% {\metacolor{\mathcal{T}_{\Ctx}( #1 , #2 )}}
\newcommand*\TTA[1]{\metacolor{\llbracket #1 \rrbracket_{\textsc t}}} % {\metacolor{\mathcal{T}_{\T} (#1)}}
\newcommand*\TS[2]{\metacolor{\llbracket #1, #2 \rrbracket_{\textsc s}}} % {\metacolor{\mathcal{T}_{\S} (#1, #2)}}
\newcommand*\TSA[1]{\metacolor{\llbracket #1 \rrbracket_{\textsc s}}}% {\metacolor{\mathcal{T}_{\S} (#1)}}
\newcommand*\TCA[1]{\metacolor{\llbracket #1\rrbracket_{\textsc c}}}
%{\metacolor{\mathcal{T}_{\C} (#1)}}
\newcommand*\TEA[1]{\metacolor{\llbracket #1 \rrbracket_{\textsc e}}}% {\metacolor{\mathcal{T}_{\E} (#1)}}

\newcommand*\Cell[2][{}]{\targettypecolor{\Keyword{Cell}^{\priocolor{#1}}} (#2)}

% Predicates on types
\newcommand*\IsBorrow[1]{\Keyword{Borrow}~#1}
\newcommand*\IsOwning[1]{\Keyword{Owning}~#1}

% General math
\newcommand*\Nat{\mathbf{N}}

% Target types
\newcommand*\HasTargetType[4][{}]{#2 \vdash^{#1} {#3}:{#4} }

% Notes
\newcommand{\note}[2]{{\color{red}[#1:} #2{\color{red}]}}
\renewcommand{\vv}[1]{\note{vv}{#1}} % clashes with \vv from newtxmath, loaded automatically by acmart
\newcommand{\pt}[1]{\note{pt}{#1}}
\newcommand{\js}[1]{\note{js}{#1}}

%%% Local Variables:
%%% mode: latex
%%% TeX-master: "main"
%%% End:

% Macro for easy \inferrule definitions.
%
%   \DefRule {⟨name⟩} {⟨premises⟩} {⟨conclusion⟩}
%
% defines a new macro \⟨prefix⟩⟨name⟩ where ⟨prefix⟩ is the definition of
% \RuleCmdPrefix. The macro contents are equivalent to
%
%   \inferrule[\RuleNamePrefix]{⟨premises⟩}{⟨conclusion⟩}
% 
% == Usage example ==
%
% Define for each group of rules the prefixes for names and commands:
%
%   \def\RuleNamePrefix{PrintedNamePrefix}
%   \def\RuleCmdPrefix{RulePrefixExample}
%
% Define rules:
%
%   \DefRule {MyRule} {...} {...}
%
% And use in the document:
%
%   \RulePrefixExampleMyRule
%
% which produces the \inferrule with the name "PrintedNamePrefixMyRule"
\NewDocumentCommand \DefRule {s o m m m} {
  \expandafter \DefRuleHelper \expandafter {\RuleNamePrefix} {#1} {#2} {#3} {#4} {#5}
}

\newcommand*\DefRuleHelper[6]{
  \IfBooleanTF {#2}
    {
      \expandafter \newcommand \csname \RuleCmdPrefix #4\endcsname
        {
          \inferrule*[#3]{#5}{#6}
        }
    }
    {
      \expandafter \newcommand \csname \RuleCmdPrefix #4\endcsname
        {
          \inferrule[#1#4]{#5}{#6} \IfValueT{#3}{~#3}
        }
    }
}

% \input{rules/session-type-formation}
% \input{rules/type-formation}

\input{rules/bounded-session-types}
\input{rules/congruence}
% Constant Typing

\def\RuleCmdPrefix{RuleCT}
\def\RuleNamePrefix{CT-}

% CT-Unit
\DefRule {Unit} {} {
  \HasCType\CUnit\BUnit
}

% CT-Fork
\DefRule {Fork} {} {
  \HasCType\CFork{
    \TArr{\MobMobile}{\DirUnord,\EffPure}{
      (\TArr{\MobMobile}{\DirUnord,\EffImpure}\BUnit\BUnit)
    }{\BUnit}
  }
}

% CT-New
\DefRule {New} {
  \SNew \S
} {
  \HasCType
    {\CNew_{\S}}
    {\TArr{\MobMobile}{\DirUnord,\EffPure}
      {\BUnit}
      {\TPairUnord
        {(\SSeq{\SAcq}{\SSeq\S\SWait})}
        {(\SSeq{\SAcq}{\SSeq{\Dual~\S}{\STerm}})}
      }
    }
}

% CT-Drop
\DefRule {Drop} {} {
  \HasCType
  {\CDrop}{\TArr{\MobMobile}{\DirUnord,\EffPure}{\SRet}{\BUnit}}
}

% CT-Acquire
\DefRule {Acquire} {} {
  \HasCType
  {\CAcq}{\TArr{\MobMobile}{\DirUnord,\EffPure}{\SSeq\SAcq\S}{\S}}
}

% CT-LSplit
\DefRule {LSplit} {
  \neg \Skips {\S[1]}
} {
  \HasCType
    {\CSplit_{\S[1],\S[2]}}
    {\TArr{\MobMobile}{\DirUnord,\EffPure}{\SSeq {\S[1]}{\S[2]}}{\TPairOrd{\S[1]}{\S[2]}}}
}

% CT-RSplit
\DefRule {RSplit} {} {
  \HasCType
  {\CSSplit_{\S[1],\S[2]}}
  {\TArr{\MobMobile}{\DirUnord,
    \EffPure}
    {\SSeq {\S[1]}{\S[2]}}
    {\TPairUnord{(\SSeq{\S[1]}{\SRet})}{(\SSeq{\SAcq}{\S[2]})}}}
}

% CT-Wait
\DefRule {Wait} {} {
  \HasCType\CWait{\TArr{\MobMobile}{\DirUnord,\EffImpure}{\SWait}\BUnit}
}

% CT-Term
\DefRule {Term} {}{
  \HasCType\CTerm{\TArr{\MobMobile}{\DirUnord,\EffImpure}{\STerm}\BUnit}
}

% CT-Send
\DefRule {Send} {
  \Mobile \T
}{
  \HasCType
    {\CSend_{\T}}
    {\TArr{\MobMobile}{\DirUnord,\EffImpure}{\TPairUnord{\T}{\SSend\T}}\BUnit}
}

% CT-Recv
\DefRule {Recv} {
  \Mobile \T
}{
  \HasCType
    {\CRecv_{\T}}
    {\TArr{\MobMobile}{\DirUnord,\EffImpure}{\SRecv\T}\T}
}

% CT-Select
\DefRule {Select} {
  k \in L
} {
  \HasCType{\CSelect{k}}{
    \TArr{\MobMobile}{\DirUnord,\EffImpure}
      {\SChoice{\ell:\S[\ell]}{\ell\in L}}
      {\S[k]}
  }
}

% \DefRule {Select} {} {
%   \HasCType{\CSelect{k}}{\TArr{\MobMobile}{}{\SChoice{k:\SSkip}{}}\BUnit}
% }

% CT-Branch
\DefRule {Branch} {} {
  \HasCType\CBranch{
    \TArr{\MobMobile}{\DirUnord,\EffImpure}
      {\SBranch{\ell:\S[\ell]}{\ell\in L}}
      {\TVariant{\ell:\S[\ell]}}
      %{\Sigma_k \S[k]}
  }
}

% CT-Discard
\DefRule {Discard} {} {
  \HasCType\CDiscard{
    \TArr{\MobMobile}{\DirUnord,\EffPure}
      {\SSkip}
      {\BUnit}
  }
}

%%%%%%%%%%%%%%%%%%%%%%%%%%%%%%%%%%%%%%%%%%%%%%%%%%%%%%%%%%%%

% \DefRule {New} {} {
%   \HasCType {\CNew_{\S}}
%     {\TArr{\MobMobile}{\DirUnord,\EffPure}{\BUnit}{\TPairUnord{\S^\blacksquare}{{\SDual\S}^\blacksquare}}}
% }

% \DefRule {Split} {} {
%   \HasCType{\CSplit_{\S[1],\S[2]}}{
%     \TArr {\MobMobile} {\DirUnord,\EffPure}
%       {{(\SSeq {\S[1]}{\S[2]})}^{\K}}
%       {\TPairOrd{\S[1]^{\S[2]\fatsemi\K}}{\S[2]^{\K}}}
%   }
% }

%%% Local Variables:
%%% mode: latex
%%% TeX-master: "../main"
%%% End:

\def\RuleNamePrefix{CS-}
\def\RuleCmdPrefix{RuleCS}

% CS-Pair
\DefRule {Pair} {} {
  \AT \TPair \ATU \CEq \AT* \TPair \ATU*
  \csred
  \left\{
    \AT \CEq \AT*\!,
    \ATU \CEq \ATU*
  \right\}
}

\DefRule {Sum} {} {
  \TSum {\AT} {\ATU} \CEq \TSum {\AT*} {\ATU*}
  \csred
  \left\{
    \AT \CEq \AT*\!,
    \ATU \CEq \ATU*
  \right\}
}

\DefRule {Unit} {} {
  \BUnit \CEq \BUnit
  \csred
  \emptyset
}

\DefRule {Arrow} {} {
  \TArr{\Mob}{\Dir,\Eff}{\AT}{\ATU}
  \CEq
  \TArr{\Mob}{\Dir,\Eff}{\AT*}{\ATU*}
  \csred
  \{
    \AT* \CEq \AT\!,
    \ATU \CEq \ATU*
  \}
}

\DefRule {Ret} {} {
  \CBounded \SRet \csred \emptyset
}

\DefRule {Close} {} {
  \CBounded \STerm \csred \emptyset
}

\DefRule {Wait} {} {
  \CBounded \SWait \csred \emptyset
}

\DefRule {SeqI} {
  \Skips \AS[2]
}{
  \CBounded\ (\SSeq{\AS[1]}{\AS[2]}) \csred \{ \CBounded \AS[1] \}
}

\DefRule {SeqII} {
  \neg \Skips \AS[2]
}{
  \CBounded\ (\SSeq{\AS[1]}{\AS[2]}) \csred \{ \CBounded \AS[2] \}
}

\DefRule {MobileUnit} {} {
  \CMobile \BUnit \csred \emptyset
}

\DefRule {Acq} {} {
  \CMobile\ (\SSeq\SAcq\AS) \csred \{ \CBounded \AS \}
}

\DefRule {MobileArrow} {} {
  \CMobile\ (\TArr{\MobMobile}{\Dir,\Eff}{\AT}{\ATU}) \csred \emptyset
}

\DefRule {MobilePair} {} {
  \CMobile\ (\AT \TPair \ATU) \csred \{ \CMobile\AT, \CMobile\ATU \}
}

\DefRule {MobileVariant} {} {
  \CMobile \TVariant{\ell:\AT[\ell]} \csred \{ \CMobile\AT[\ell] \mid \ell\in L \}
}

\DefRule {Context} {
  \AT \CEq \ATU \csred \CSet*
}{
  \CSet \cup \{ \AT\CEq\ATU \}
  \csred
  \CSet \cup \CSet*
}

\input{rules/context-formation}
% Predicate for session types creatable by `new`.

\def\RuleCmdPrefix{RuleSNew}
\def\RuleNamePrefix{TN-}

% TN-Skip
\DefRule {Skip} {} {
  \SNew \SSkip
}

% TN-Seq
\DefRule {Seq} {
  \SNew {\S[1]} \\
  \SNew {\S[2]}
} {
  \SNew {\SSeq{\S[1]}{\S[2]}}
}

% TN-Rec
\DefRule {Rec} {
  \SNew {\S}
} {
  \SNew {\SRec X \S}
}

% TN-Var
\DefRule {Var} {} {
  \SNew {X}
}

% TN-Branch
\DefRule {Branch} {
  \forall \ell.\ \SNew {\S[\ell]}
} {
  \SNew {\SBranch{\ell:\S[\ell]}{\ell\in L}}
}

% TN-Choice
\DefRule {Choice} {
  \forall \ell.\ \SNew {\S[\ell]}
} {
  \SNew {\SChoice{\ell:\S[\ell]}{\ell\in L}}
}

% TN-Send
\DefRule {Send} {} {
  \SNew {\SSend \T}
}

% TN-Recv
\DefRule {Recv} {} {
  \SNew {\SRecv \T}
}

\input{rules/ord-types}
% Process typing

\newcommand*\RuleTExpr{
  \inferrule[TP-Expr]{
    \HasType[\EffImpure]
      \Ctx
      \E
      \BUnit
  }{
    \ProcHasType
    \Ctx
    {\PExp{\E}}
  }
}

\newcommand*\RuleTPar{
  \inferrule[TP-Par]{
    \ProcHasType{\Ctx[1]}{\Proc[1]} \\
    \ProcHasType{\Ctx[2]}{\Proc[2]}
  }{
    \ProcHasType
    {\Ctx[1] \CtxPar \Ctx[2]}
    {\PPar{\Proc[1]}{\Proc[2]}}
  }
}

\newcommand*\RuleTRes{
  \inferrule[TP-Res]{
    \ProcHasType
      {\Ctx \CtxPar \Ctx[1] \CtxPar \Ctx[2]}
      \Proc  \\
    \BindGivesCtx \S {\BindGroup[1]} {\Ctx[1]} \\
    \BindGivesCtx {\SDual\S} {\BindGroup[2]} {\Ctx[2]}
  }{
    \ProcHasType
    {\Ctx}
    {\PNew{\BindGroup[1]}{\BindGroup[2]}\Proc}
  }
}

\newcommand*\RuleBindEmp{
  \inferrule[B-Emp]{}{
    \BindGivesCtx \SSkip \BEmp \CtxEmpty
  }
}

\newcommand*\RuleBindSeq{
  \inferrule[B-Seq]{
    \BindGivesCtx {\S[1]} \EVar {\Ctx[1]} \\
    \BindGivesCtx {\S[2]} \Bind {\Ctx[2]}
  }{
    \BindGivesCtx
      {\SSeq{\S[1]}{\S[2]}}
      {\BSeq{\EVar}{\Bind}}
      {\Ctx[1] \CtxSeq \Ctx[2]}
  }
}

\newcommand*\RuleBindSep{
  \inferrule[B-Sep]{
    \BindGivesCtx {\SSeq{\S[1]}{\SRet}} {\Bind[1]} {\Ctx[1]} \\
    \BindGivesCtx {\SSeq{\SAcq}{\S[2]}} {\BindGroup[2]} {\Ctx[2]}
  }{
    \BindGivesCtx
      {\SSeq{\S[1]}{\S[2]}}
      {\BWithSep{\Bind[1]}{\BindGroup[2]}}
      {\Ctx[1] \CtxPar \Ctx[2]}
  }
}

\newcommand*\RuleBindVar{
  \inferrule[B-Var]{}{
    \BindGivesCtx \S \EVar {\CtxVar\EVar\S}
  }
}

%%% Local Variables:
%%% mode: latex
%%% TeX-master: "../main"
%%% End:

\input{rules/reduction}
\input{rules/session-subtyping}
% Congruence of translated processes

\def\RuleNamePrefix{TrC-}
\def\RuleCmdPrefix{RuleTransCong}

% TrC-Refl
\DefRule {Refl} {} {
  \TProc \Congruent \TProc
}

% TrC-ParComm
\DefRule {ParComm} {} {
  \TPPar{\TProc[1]}{\TProc[2]}
    \Congruent \TPPar{\TProc[2]}{\TProc[1]}
}

% TrC-ParAssoc
\DefRule {ParAssoc} {} {
  \TPPar{\TProc[1]}{(\TPPar{\TProc[2]}{\TProc[3]})}
    \Congruent \TPPar{(\TPPar {\TProc[1]}{\TProc[2]})}{\TProc[3]}
}

% TrC-Trans
\DefRule {Trans} {
  \TProc[1] \Congruent \TProc[2] \\
  \TProc[2] \Congruent \TProc[3]
}{
  \TProc[1] \Congruent \TProc[3]
}

% TrC-ParUnit
\DefRule {ParUnit} {} {
  \TPPar{\TPExp{\CUnit}}{\TProc} \Congruent \TProc
}

% TrC-ResSwap
\DefRule {ResSwap} {} {
  \TPNew[xy] \TProc
    \Congruent \TPNew[yx] \TProc
}

% TrC-ResComm
\DefRule {ResComm} {
  \{\TEVar[1],\TEVary[1]\} \cap \{\TEVar[2],\TEVary[2]\} = \emptyset
}{
  \TPNew[x_1y_1] {\TPNew[x_2y_2] \TProc}
  \Congruent
  \TPNew[x_2y_2] {\TPNew[x_1y_1] \TProc} 
}

% TrC-FlagComm
\DefRule {FlagComm} {
  \TEVarz[1] \neq \TEVarz[2]
}{
  \TPSync [{\TEVarz[1]}] {\Flag[1]} \TPSync [{\TEVarz[2]}] {\Flag[2]} \TProc
  \Congruent
  \TPSync [{\TEVarz[2]}] {\Flag[2]} \TPSync [{\TEVarz[1]}] {\Flag[1]} \TProc 
}

% TrC-ResFlagComm
\DefRule {ResFlagComm} {
  \TEVarz \not\in \{\TEVar,\TEVary\}
}{
  \TPNew {\TPSync \Flag \TProc}
  \Congruent
  \TPSync \Flag {\TPNew \TProc} 
}


% TrC-ResExtend
\DefRule {ResExtend} {
  \TEVar,\TEVary \not\in \Fv(\TProc[1])
}{
  \TPPar{\TProc[1]}{\TPNew{\TProc[2]}}
  \Congruent
  {\TPNew{(\TPPar{\TProc[1]}{\TProc[2]})}}
}

% TrC-FlagExtend
\DefRule {FlagExtend} {
  \TEVarz \not\in \Fv(\TProc[1])
}{
  \TPPar{\TProc[1]}{\TPSync \Flag {\TProc[2]}}
  \Congruent
  {\TPSync \Flag {(\TPPar{\TProc[1]}{\TProc[2]})}}
}


%%% Local Variables:
%%% mode: latex
%%% TeX-master: "../main"
%%% End:

%%%%%%%%%%%%%%%%%%%%%%%%%%%%%%%%%%%%%%%%%%%%%%%%%%%%%%%%%%%%%%%%%%%%%%%%%%%%%%%
% Reduction rules

\def\RuleCmdPrefix{RuleUProcRed}
\def\RuleNamePrefix{RU-}

% RU-Exp
\DefRule {Exp} {
  \E[1]  \Reduces \E[2]
}{
  \TFCapp{\E[1]} \Reduces \TFCapp{\E[2]}
}

% RU-Fork
\DefRule {Fork} {} {
  \TFCapp{\EApp \CFork \V}
  \Reduces
  \TPPar{\TFCapp\CUnit}{\TPExp{\EApp\V\CUnit}}
}

% RU-New
\DefRule {New} {} {
  \TFCapp{\EApp \CNew \CUnit}
  \Reduces
  \TPNew {
    \TPSync [z_1] \FlagDropped
    \TPSync [z_2] \FlagDropped
    (
      \TFCapp{\EPair
        {(\TEChan [\TEVarz[1]] \TEVar)}
        {(\TEChan [\TEVarz[2]] \TEVary)}
      }
    )
  }
}

% RU-LSplit
\DefRule {LSplit} {} {
    \TPNew {(
      \TPPar {\TFCapp{\EApp \CSplit {\ChanPat\TEVar}}} \TProc
    )}
    \Reduces
    \TPNew {(
      \TPPar
        {\TFCapp{\EPair{\ChanPat\TEVar}{\ChanPat\TEVar}}}
        \TProc
    )}
}

% RU-RSplit
\DefRule {RSplit} {} {
  \TPNew {(
    \TPPar {\TFCapp{\EApp \CSSplit {(\TEChan [e_1] \TEVar [e_2])}}} \TProc
  )}
  \Reduces
  \TPNew {
    \TPSync \FlagInit (
      \TPPar {
        \TFCapp{\EPair
          {(\TEChan [e_1]     \TEVar [\TEVarz])}
          {(\TEChan [\TEVarz] \TEVar [e_2])}
        }
      } \TProc
    )
  }
}

% RU-Cleanup
\DefRule {Cleanup} {} {
  \TPSync \FlagAcquired \TProc
  \Reduces
  \Subs \TProc \TEVarz \CUnit
}

% RU-Drop
\DefRule {Drop} {} {
    \TPNew {\TPSync \FlagInit (
      \TPPar {\TFCapp{\EApp \CDrop {(\TEChan [\E] \TEVar [\TEVarz])}}} \TProc
    )}
    \Reduces 
    \TPNew {\TPSync \FlagDropped (
      \TPPar {\TFCapp{\CUnit}} \TProc
    )}
}

% RU-Acquire
\DefRule {Acquire} {} {
    \TPNew {
      \TPSync \FlagDropped 
        (\TPPar {\TFCapp{\EApp \CAcq {(\TEChan [\TEVarz] \TEVar [e])}}} \TProc)
    }
    \Reduces 
    \TPNew {
      \TPSync \FlagAcquired 
        (\TPPar {\TFCapp{\TEChan \TEVar[e]}} \TProc)
    }
}

% RU-Close
\DefRule {Close} {} {
    \TPNew {(
      \TPPar {\TFCapp[1]{\EApp\CTerm {\ChanPat[1]\TEVar}}}
             {\TFCapp[2]{\EApp\CWait {\ChanPat[2]\TEVary}}}
    )}
    \Reduces 
    \TPPar {\TFCapp[1]{\CUnit}} {\TFCapp[2]{\CUnit}}
}

% RU-Com
\DefRule {Com} {} {
  \TPNew {(
    \TPPar {\TFCapp[1]{\EApp {\CSend} {(\EPair \TV {\ChanPat[1]\TEVar})}}} {
    \TPPar {\TFCapp[2]{\EApp {\CRecv} {\ChanPat[1]\TEVary}}} {\TProc}}
  )}
  \Reduces
  \TPNew {(
    \TPPar {\TFCapp[1]{\CUnit}} {
    \TPPar {\TFCapp[2]{\TV}} {\TProc}}
  )}
}

% RU-Choice
\DefRule {Choice} {} {
  \TPNew {(
    \TPPar {\TFCapp[1]{\EApp {\CSelect k} {\ChanPat[1] \TEVar}}} {
    \TPPar {\TFCapp[2]{\EApp {\CBranch}   {\ChanPat[2] \TEVary}}} {\TProc}}
  )}
  \Reduces\\\\
  \TPNew {(
    \TPPar {\TFCapp[1]{\ChanPat[1] \TEVar}} {
    \TPPar {\TFCapp[2]{\EInj k {\ChanPat[2] \TEVary}}} {\TProc}}
  )}
}

% RU-Struct
\DefRule {Struct} {
  \TProc[1] \equiv \TProc*[1] \\
  \TProc*[1] \Reduces \TProc*[2] \\
  \TProc*[2] \equiv \TProc[2]
}{
  \TProc[1] \Reduces \TProc[2]
}

% RU-Par
\DefRule {Par} {
  \TProc[1] \Reduces \TProc[2]
}{
  \TPPar {\TProc[1]}{\TProcQ} \Reduces \TPPar {\TProc[2]}{\TProcQ}
}

% RU-Bind
\DefRule {Bind} {
  \TProc[1] \Reduces \TProc[2]
}{
  \TPNew {\TProc[1]} \Reduces \TPNew {\TProc[2]}
}

% RU-Sync
\DefRule {Sync} {
  \TProc[1] \Reduces \TProc[2]
}{
  \TPSync \Flag {\TProc[1]} \Reduces \TPSync \Flag {\TProc[2]}
}


%%%%%%%%%%%%%%%%%%%%%%%%%%%%%%%%%%%%%%%%%%%%%%%%%%%%%%%%%%%%%%%%%%%%%%%%%%%%%%%
% Translation of processes

\def\RuleCmdPrefix{RuleTransProc}
\def\RuleNamePrefix{Tr-}

% Tr-Exp
\DefRule {Exp} {} {
  \PExp{\E}
  \TransProc
  \TPExp{\metacolor{\Sub{[\E*]}}}
}

% Tr-Par
\DefRule {Par} {
  \Proc[1] \TransProc \TProc*[1] \\
  \Proc[2] \TransProc \TProc*[2]
} {
  \PPar {\Proc[1]} {\Proc[2]}
  \TransProc
  \TPPar {\TProc*[1]} {\TProc*[2]}
}

% Tr-Bind
\DefRule {Bind} {
  \BindGroup[1] \TransBind[\TEChan \TEVar] {\TCC[1]}  [\Sub[1]] \\
  \BindGroup[2] \TransBind[\TEChan \TEVary] {\TCC[2]} [\Sub[2]] \\
  \Proc \TransProc[{\Sub \cup \Sub[1] \cup \Sub[2]}] \TProc* \\
} {
  \PNew {\BindGroup[1]}{\BindGroup[2]} \Proc
  \TransProc
  \TPNew {\TCC[1][\TCC[2][\TProc*]]}
}

%%%%%%%%%%%%%%%%%%%%%%%%%%%%%%%%%%%%%%%%%%%%%%%%%%%%%%%%%%%%%%%%%%%%%%%%%%%%%%%
% Translation of bind groups

\def\RuleCmdPrefix{RuleTransBind}
\def\RuleNamePrefix{Tr-Bind}

% Tr-BindPar
\DefRule {Par} {
  \BindGroup[1] \TransBind[{\TEChan[\E[1]] \TEVar [\TEVarz]}] {\TCC[1]} [\Sub[1]] \\
  \BindGroup[2] \TransBind[{\TEChan[\TEVarz] \TEVar [\E[1]]}] {\TCC[2]} [\Sub[2]] \\
  \BindFlag{\BindGroup[1]} = \Flag
}{
  \BPar {\BindGroup[1]} {\BindGroup[2]} \TransBind[{\TEChan[\E[1]] \TEVar [\E[2]]}]
    {\TPSync \Flag \targetproccolor{(\TCC[1][\TCC[2]])}}
    [\Sub[1] \cup \Sub[2]]
}

% Tr-BindChannels
\DefRule {Channels} {
  \Bind \Downarrow \Bind* \\
  \Bind* \TransBind \Sub
}{
  \Bind \TransBind \TCCEmpty [\Sub]
}

% Tr-BindSeq
\DefRule {Seq} {
  \Bind[1] \TransBind[\TEChan[\E[1]]{\TEVar}] {\Sub[1]} \\
  \Bind[2] \TransBind[\TEChan{\TEVar}[\E[2]]] {\Sub[2]} \\
}{
  \BSeq{\Bind[1]}{\Bind[2]} \TransBind[\TEChan[\E[1]]{\TEVar}[\E[2]]] {\Sub[1] \cup \Sub[2]}
}

% Tr-BindVar
\DefRule {Var} {} {
  \EVar \TransBind \{ \EVar\mapsto\E \}
}

% Tr-BindNil
\DefRule {Nil} {} {
  \BEmp \TransBind \emptyset
}

% Type Context Splitting


\newcommand*\RuleCSEmpty{
  \inferrule[CS-Empty]{}{
    \CtxSplit\CtxEmpty\CtxEmpty\CtxEmpty\D
  }
}

\newcommand*\RuleCSSplit{
  \inferrule[CS-Split]{
    \CtxSplit
      {\Ctx}
      {\Ctx_1}
      {\Ctx_2}
      {\D\cup\D'}
    \\
    \EVar \not\in \D
  }{
    \CtxSplit
      {\Ctx \CtxSeq \EVar{(\SSeq{\S_1}{\S_2})}{\D'}}
      {\Ctx[1] \CtxSeq \EVar{\S_1}{\D'}}
      {\Ctx[2] \CtxSeq \EVar{\S_2}{\D'}}
      {\D}
  }
}

\newcommand*\RuleCSLeft{
  \inferrule[CS-Left]{
    \CtxSplit
      {\Ctx}
      {\Ctx[1]}
      {\Ctx[2]}
      {\D}
    \\
    \EVar \not\in \D
  }{
    \CtxSplit
      {\Ctx \CtxSeq \EVar\T{\D'}}
      {\Ctx[1] \CtxSeq \EVar\T{\D'}}
      {\Ctx[2]}
      {\D}
  }
}

\newcommand*\RuleCSRight{
  \inferrule[CS-Right]{
    \CtxSplit
      {\Ctx}
      {\Ctx_1}
      {\Ctx_2}
      {\D\cup\D'}
  }{
    \CtxSplit
      {\CtxExt\Ctx\EVar\T{\D'}}
      {\CtxExtInactive{\Ctx_1}\EVar\T{\D'}}
      {\CtxExt{\Ctx_2}\EVar\T{\D'}}
      {\D}
  }
}

\newcommand*\RuleCSBoth{
  \inferrule[CS-Both]{
    \CtxSplit
      {\Ctx}
      {\Ctx_1}
      {\Ctx_2}
      {\D}
    \\
    \Unr \T
  }{
    \CtxSplit
      {\CtxExt\Ctx\EVar\T{\emptyset}}
      {\CtxExt{\Ctx_1}\EVar\T{\emptyset}}
      {\CtxExt{\Ctx_2}\EVar\T{\emptyset}}
      {\D}
  }
}

% Type Concatenation

\def\RuleNamePrefix{TC-}
\def\RuleCmdPrefix{TypeConcatenation}

\DefRule {Skip} { } {
    \ConcatType {\T} {\SSkip} = \T
}

\DefRule {SemiDist} {
    V \neq \SSkip
}{
    \ConcatType {(\SSeq {\T}{\TU})} {V} = \SSeq {\T} {(\ConcatType {\TU} {V})}
}

\DefRule {NonSemi} {
    \T \neq \SSeq {V} {W} \\
    \TU \neq \SSkip
}{
    \ConcatType {\T} {\TU} = \SSeq {\T} {\TU}
}

% Type Normalisation

\def\RuleNamePrefix{TN-}
\def\RuleCmdPrefix{TypeNormalisation}

\DefRule {SemiSkips} {
    {\T} \SEq \SSkip \\
    \NormaliseType {\TU} {\TU'}
}{
    \NormaliseType {\SSeq \T \TU} {\TU'}
}

\DefRule {Semi} {
    \neg ({\T} \SEq \SSkip) \\
    \NormaliseType {\T} {\T'}
}{
    \NormaliseType {\SSeq {\T} \TU} {\ConcatType {\T'} {\TU}}
}

\DefRule {Mu} {
    \NormaliseType \S {\S'}
}{
    \NormaliseType {\SRec \SVar \S} {{\S'} [{\SRec \SVar {\S'}} / \SVar]}
}

\DefRule {ChoiceSemi} { }{
    \NormaliseType  {\SSeq {\SChoice {\ell:\S[\ell]} {\ell\in L}} {\T}} {\SChoice {\ell: \SSeq {\S[\ell]} {\T}} {\ell\in L}}
}

\DefRule {BranchSemi} { }{
    \NormaliseType  {\SSeq {\SBranch {\ell:\S[\ell]} {\ell\in L}} {\T}} {\SBranch {\ell: \SSeq {\S[\ell]} {\T}} {\ell\in L}}
}

\DefRule {Identity} {
    \T \neq (\SSeq {\TU} {\T}, \SRec {\SVar} {\S}, \SChoice {\ell:\S[\ell]} {\ell\in L})
}{
    \NormaliseType {\T} {\T}
}

\input{rules/type-subtyping}
% Algorithmic Expression Typing

\def\RuleNamePrefix{A-}
\def\RuleCmdPrefix{RuleAlg}

% A-Const
\DefRule [(\C \not\in \{\CSplit,\CSSplit\})] {Const} {
  \CtxEmpty \SubCtx \Ctx \\
  \HasCType \C \T
}{
  \InferType \C \T \EffPure [\CSEmpty]
}

% A-LSplit
\DefRule {LSplit} {
  \CtxEmpty \SubCtx \Ctx \\
  \TVar\ \text{fresh}
}{
  \InferType {\CSplit_{\S}}
  {\TArr{\MobMobile}{\DirUnord,\EffPure} {\SSeq\S\TVar} {\TPairOrd{\S}{\TVar}}}
    \EffPure
    [\CSEmpty]
}

% A-RSplit
\DefRule {RSplit} {
  \CtxEmpty \SubCtx \Ctx \\
  \TVar\ \text{fresh}
}{
  \InferType {\CSSplit_{\S}}
  {\TArr{\MobMobile}{\DirUnord,\EffPure} {\SSeq\S\TVar} {\TPairUnord{\SSeq{\S}{\SRet}}{\SSeq\SAcq\TVar}}}
    \EffPure
    [\CSEmpty]
}

% A-Var
\DefRule {Var} {
  \CtxVar \EVar \Ctx(\EVar) \SubCtx \Ctx
} {
  \InferType \EVar {\Ctx(\EVar)} \EffPure [\CSEmpty]
}

% A-App
\DefRule {App} {
  \CtxAlgJoinCheck{\E[1]}{\E[2]} \\
  \Dir = \DirLeft \implies \Eff[1] = \EffPure \\
  \Dir = \DirRight \implies \Eff[2] = \EffPure \\\\
  \InferType [\CtxRestrict{\E[1]}] {\E[1]} {\TArr{\Mob}{\Dir,\Eff[3]} \T \TU} {\Eff[1]} [\CSet[1]] \\
  \CheckType [\CtxRestrict{\E[2]}] {\E[2]} {\T} {\Eff[2]} [\CSet[2]]
}{
  \InferType {\E[1]\E[2]} {\TU} {\Eff[1] \EffJoin \Eff[2] \EffJoin \Eff[3]} [\CSet[1] \CSJoin \CSet[2]]
}

% A-Check
\DefRule {Check} {
  \InferType \EVar {\T[2]} {\Eff} \\
}{
  \CheckType \EVar {\T[1]} {\Eff} [\CSet \CSJoin \T[2] \CSub \T[1]]
}

% A-Abs
\DefRule {Abs} {
  \CheckType[\CtxVar\EVar{\T[1]} \CtxJoin \Ctx]
    \E
    {\T[2]}
    {\Eff*} \\
  \Eff* \le \Eff \\
  \Dir = \DirNone \implies \Unr \Ctx
}{
  \CheckType
    {\EAbs\EVar\E}
    {\TArr{\Mob}{\Dir,\Eff}{\T[1]}{\T[2]}}
    {\EffPure}
}

% A-AbsRec
\DefRule {AbsRec} {
  \CheckType[
    \CtxVar f {\TArr{\Mob}{\DirNone,\Eff} \T \TU}
      \CtxPar \CtxVar\EVar\T
      \CtxPar \Ctx
  ]
    \E
    \TU
    {\Eff*} \\
  \Eff* \le \Eff \\
  \Unr\Ctx \\
}{
  \CheckType
    {\EMu {\EAbs\EVar\E}}
    {\TArr{\Mob}{\DirNone,\Eff} \T \TU}
    {\EffPure}
}

% A-Pair
\DefRule {Pair} {
  \CtxAlgJoinCheck{\E[1]}{\E[2]} \\
  \Dir = \DirLeft \implies \Eff[2] = \EffPure \\
  \Dir = \DirRight \implies \Eff[1] = \EffPure \\\\
  \CheckType[\CtxRestrict{\E[1]}]
    {\E[1]}
    {\T}
    {\Eff[1]}
    [\CSet[1]]
  \\
  \CheckType[\CtxRestrict{\E[2]}]
    {\E[2]}
    {\TU}
    {\Eff[2]}
    [\CSet[2]]
} {
  \CheckType
    {\EPair {\E[1]} {\E[2]}}
    {\T \TPair \TU}
    {\Eff}
    [\CSet[1] \CSJoin \CSet[2]]
}

% \newcommand*\RuleTPairUnord{
%   \inferrule[T-PairUnord]{
%     \HasType[\Eff[1]]
%       {\Ctx_1}
%       {\E[1]}
%       {\T[1]} \\
%     \HasType[\Eff[2]]
%       {\Ctx_2}
%       {\E[2]}
%       {\T[2]}
%   }{
%     \HasType[\Eff[1] \EffJoin \Eff[2]]
%       {\Ctx[1] \CtxPar  \Ctx[2]}
%       {\EPair{\E[1]}{\E[2]}}
%       {\TPairUnord{\T[1]}{\T[2]}}
%   }
% }
% 
% \newcommand*\RuleTPairOrd{
%   \inferrule[T-PairOrd]{
%     % \Unr{\Ctx_1} \vee \Unr {\T[1]} \vee \Eff[2] = \EffPure \\
%     \HasType[\Eff[1]]
%       {\Ctx_1}
%       {\E[1]}
%       {\T[1]} \\
%     \HasType[\EffPure]
%       {\Ctx_2}
%       {\E[2]}
%       {\T[2]}
%   }{
%     \HasType[\Eff[1]]
%       {\Ctx[1] \CtxSeq  \Ctx[2]}
%       {\EPair{\E[1]}{\E[2]}}
%       {\TPairOrd{\T[1]}{\T[2]}}
%   }
% }

% A-LetUnit
\DefRule {LetUnit} {
  \CtxAlgSeqCheck{\E[1]}{\E[2]} \\
  \CheckType[\CtxRestrict{\E[1]}]
    {\E[1]}
    {\BUnit}
    {\Eff[1]}
    [\CSet[1]]
  \\
  \InferType[\CtxRestrict{\E[2]}]
    {\E[2]}
    {\T}
    {\Eff[2]}
    [\CSet[2]]
}{
  \InferType
    {\ESeq{\E[1]}{\E[2]}}
    {\T}
    {\Eff[1] \EffJoin \Eff[2]}
    [\CSet[1] \CSJoin \CSet[2]]
}

% \newcommand*\RuleTLet{
%   \inferrule[T-Let]{
%     \HasType[\Eff[1]]
%       {\Ctx[1]}
%       {\E[1]}
%       {\T[1]} \\
%     \HasType[\Eff[2]]
%       {\CtxVar x {{\T[1]}} \CtxPatOp \Ctx[2]}
%       {\E[2]}
%       {\T[2]}
%   }{
%     \HasType[\Eff[1] \EffJoin \Eff[2]]
%       {\CtxPat 1 2}
%       {\ELet x{\E[1]}{\E[2]}}
%       {\T[2]}
%   }
% }

% A-LetPair
\DefRule {LetPair} {
  \CtxAlgSeqCheck{\E[1]}{\E[2]} \\
  \InferType[\CtxRestrict{\E[1]}]
    {\E[1]}
    {\T[1] \TPair \T[2]}
    {\Eff[1]}
    [\CSet[1]]
  \\
  \InferType[(\CtxVar x {\T[1]} \CtxJoin \CtxVar y {\T[2]}) \CtxSeq \CtxRestrict{\E[2]}]
    {\E[2]}
    {\TU}
    {\Eff[2]}
    [\CSet[2]]
}{
  \InferType
     {\ELetPair xy{\E[1]}{\E[2]}}
     {\TU}
     {\Eff[1] \EffJoin \Eff[2]}
     [\CSet[1] \CSJoin \CSet[2]]
}

% \newcommand*\RuleTLetUnord{
%   \inferrule[T-LetUnord]{
%     \HasType[\Eff[1]]
%       {\Ctx[1]}
%       {\E[1]}
%       {\TPairUnord {\T[1]}{\T[2]}} \\
%     \HasType[\Eff[2]]
%       {\CtxVar x{\T[1]} \CtxPar \CtxVar y{\T[2]} \CtxPatOp \Ctx[2]}
%       {\E[2]}
%       {\T[3]}
%   }{
%     \HasType[\Eff[1] \EffJoin \Eff[2]]
%       {\CtxPat 1 2}
%       {\ELetPair xy{\E[1]}{\E[2]}}
%       {\T[3]}
%   }
% }
% 
% \newcommand*\RuleTLetOrd{
%   \inferrule[T-LetOrd]{
%     \HasType[\Eff[1]]
%       {\Ctx[1]}
%       {\E[1]}
%       {\TPairOrd {\T[1]}{\T[2]}} \\
%     \HasType[\Eff[2]]
%       {\CtxVar x{\T[1]} \CtxSeq \CtxVar y{\T[2]} \CtxPatOp \Ctx[2]}
%       {\E[2]}
%       {\T[3]}
%   }{
%     \HasType[\Eff[1] \EffJoin \Eff[2]]
%       {\CtxPat 1 2}
%       {\ELetPair xy{\E[1]}{\E[2]}}
%       {\T[3]}
%   }
% }
% 
% \newcommand*\RuleTInLeft{
%   \inferrule[T-Inj]{
%     \HasType[\Eff]
%       {\Ctx}
%       {\E}
%       {\T[i]}
%   }{
%     \HasType[\Eff]
%       {\Ctx}
%       {\EInj{i} \E}
%       {\TSum{\T[1]}{\T[2]}}
%   }
% }
% 
% \newcommand*\RuleTCaseSum{
%   \inferrule[T-CaseSum]{
%     \HasType[\Eff]
%       {\Ctx[1]}
%       {\E}
%       {\TSum{\T[1]}{\T[2]}}
%     \\
%     \HasType[\Eff[1]]
%       {\CtxVar{\EVar_1}{\T[1]} \CtxPatOp \Ctx[2]}
%       {\E[1]}
%       {\T}
%     \\
%     \HasType[\Eff[2]]
%       {\CtxVar{\EVar_2}{\T[2]} \CtxPatOp \Ctx[2]}
%       {\E[2]}
%       {\T}
%   }{
%     \HasType[\Eff \EffJoin \Eff[1] \EffJoin \Eff[2]]
%       {\CtxPat 1 2}
%       {\ECase\E{x_1}{\E[1]}{x_2}{\E[2]}}
%       {\T}
%   }
% }
% 
% \newcommand*\RuleTWeakC{
%   \inferrule[T-WeakC]{
%     \Unr\Ctx_2 \\
%     \HasType[\Eff]
%       {\Ctx[1] \CtxSeq \Ctx[2]}
%       {\E}
%       {\T}
%   }{
%     \HasType[\Eff]
%       {\Ctx[1]}
%       {\E}
%       {\T}
%   }
% }
% 
% \newcommand*\RuleTWeakE{
%   \inferrule[T-WeakE]{
%     \HasType[\EffPure]
%       {\Ctx}
%       {\E}
%       {\T}
%   }{
%     \HasType[\EffPure]
%       {\Ctx}
%       {\E}
%       {\T}
%   }
% }
% 
% \newcommand*\RuleTWeakCE{
%   \inferrule[T-WeakCE]{
%     \Ctx[2] \preccurlyeq \Ctx[1]
%     \\
%     \Eff[1] \le \Eff[2]
%     \\
%     \HasType[\Eff[1]]
%       {\Ctx_1}
%       {\E}
%       {\T}
%   }{
%     \HasType[\Eff[2]]
%       {\Ctx_2}
%       {\E}
%       {\T}
%   }
% }
% 
% \newcommand*\RuleTMatch{
%   \inferrule[T-MatchLeft]{
%     \HasType[\Eff]
%       {\Ctx_1}
%       {\E}
%       {\SBranch{\ell:\S_\ell}{\ell\in L}}
%     \\
%     \HasType[\Eff[\ell]]
%       {\Ctx_2}
%       {\E_\ell}
%       {\TArr{\Mob}{\Eff[\ell]}{\S_\ell}{\T}}
%   }{
%     \HasType[\Eff \EffJoin \EffJoin \{ \Eff[\ell] \}_{\ell\in L}]
%       {\Ctx[1] \CtxSeq \Ctx[2]}
%       {\EMatch{\E}{\ell:\E_\ell}{\ell\in L}}
%       {\T}
%   }
% }
% 
% \newcommand*\RuleTMatchLin{
%   \inferrule[T-MatchLin]{
%     \HasType[\Eff]
%       {\Ctx_1}
%       {\E}
%       {\SBranch{\ell:\S_\ell}{\ell\in L}}
%     \\
%     \HasType[\Eff[\ell]]
%       {\Ctx_2}
%       {\E_\ell}
%       {\TArr{\Mob}{\Eff[\ell]}{\S_\ell}{\T}}
%   }{
%     \HasType[\Eff \EffJoin \EffJoin\{ \Eff[\ell]\}_{\ell\in L}]
%       {\Ctx[1] \CtxPar \Ctx[2]}
%       {\EMatch{\E}{\ell:\E_\ell}{\ell\in L}}
%       {\T}
%   }
% }
% 

%%% Local

%%% Local Variables:
%%% mode: latex
%%% TeX-master: "../main"
%%% End:

% Declarative Expression Typing

\def\RuleNamePrefix{T-}
\def\RuleCmdPrefix{RuleT}

% T-Const
\DefRule {Const} {
  \HasCType
    \C
    \T
}{
  \HasType[\EffPure]
    \CtxEmpty
    \C
    \T
}

% T-Var
\DefRule {Var} {} {
  \HasType[\EffPure]
    {\CtxVar \EVar\T}
    \EVar
    \T
}

% T-AbsUnr
\DefRule {AbsUnr} {
  \HasType
    {
        \CtxVar f {\TArr{\MobMobile}{\DirNone,\Eff}{\T}{\TU}}
          \CtxPar \CtxVar \EVar{\T}
          \CtxPar \Ctx
    }
    {\E}
    {\TU} \\
  \Unr\Ctx
}{
  \HasType[\EffPure]
    \Ctx
    {\EMu {\EAbs\EVar\E}}
    {\TArr{\MobMobile}{\DirNone,\Eff} \T \TU}
}

% T-AbsLin
\DefRule {AbsLin} {
  \HasType
    {\CtxVar\EVar{\T} \CtxPar \Ctx}
    {\E}
    {\TU} \\
  \Mob (\Ctx)
}{
  \HasType[\EffPure]
    \Ctx
    {\EAbs\EVar\E}
    {\TArr{\Mob}{\DirUnord,\Eff} \T \TU}
}

% T-AbsLeft
\DefRule {AbsLeft} {
  \HasType
    {\CtxVar{\EVar}{\T} \CtxSeq \Ctx}
    {\E}
    {\TU} \\
  \Mob (\Ctx)
}{
  \HasType[\EffPure]
    \Ctx
    {\EAbs\EVar\E}
    {\TArr{\Mob}{\DirLeft,\Eff} \T \TU}
}

% T-AbsRight
\DefRule {AbsRight} {
  \HasType
    {\Ctx \CtxSeq \CtxVar{\EVar}{\T}}
    {\E}
    {\TU} \\
  \Mob (\Ctx)
}{
  \HasType[\EffPure]
    \Ctx
    {\EAbs\EVar\E}
    {\TArr{\Mob}{\DirRight,\Eff} \T \TU}
}

% T-AppUnr
\DefRule {AppUnr} {
  \HasType
    {\Ctx[1]}
    {\E[1]}
    {\TArr{\Mob}{\DirNone,\Eff} \T \TU} \\
  \HasType
    {\Ctx[2]}
    {\E[2]}
    {\T}
}{
  \HasType
    {\Ctx[1] \CtxSeq \Ctx[2]}
    {\E[1] \E[2]}
    {\TU}
}

% T-AppLine
\DefRule {AppLin} {
  \HasType
    {\Ctx[1]}
    {\E[1]}
    {\TArr{\Mob}{\DirUnord,\Eff} \T \TU} \\
  \HasType
    {\Ctx[2]}
    {\E[2]}
    {\T}
}{
  \HasType
    {\Ctx[1] \CtxPar \Ctx[2]}
    {\E[1] \E[2]}
    {\TU}
}

% T-AppLeft
\DefRule {AppLeft} {
  \HasType[\EffPure]
    {\Ctx[1]}
    {\E[1]}
    {\TArr{\Mob}{\DirLeft,\Eff} \T \TU}
  \\
  \HasType
    {\Ctx[2]}
    {\E[2]}
    {\T}
}{
  \HasType
    {\Ctx[2] \CtxSeq \Ctx[1]}
    {\E[1] \E[2]}
    {\TU}
}

% T-AppRight
\DefRule {AppRight} {
  \HasType
    {\Ctx[1]}
    {\E[1]}
    {\TArr{\Mob}{\DirRight,\Eff} \T \TU}
    \\
  \HasType[\EffPure]
    {\Ctx[2]}
    {\E[2]}
    {\T}
}{
  \HasType
    {\Ctx[1] \CtxSeq  \Ctx[2]}
    {\E[1] \E[2]}
    {\TU}
}

% T-PairUnord
\DefRule {PairUnord} {
  \HasType
    {\Ctx[1]}
    {\E[1]}
    {\T} \\
  \HasType
    {\Ctx[2]}
    {\E[2]}
    {\TU}
}{
  \HasType
    {\Ctx[1] \CtxPar  \Ctx[2]}
    {\EPair{\E[1]}{\E[2]}}
    {\TPairUnord{\T}{\TU}}
}

% T-PairOrd
\DefRule {PairOrd} {
  % \Unr{\Ctx[1]} \vee \Unr {\T[1]} \vee \Eff[2] = \EffPure \\
  \HasType
    {\Ctx[1]}
    {\E[1]}
    {\T} \\
  \HasType[\EffPure]
    {\Ctx[2]}
    {\E[2]}
    {\TU}
}{
  \HasType
    {\Ctx[1] \CtxSeq  \Ctx[2]}
    {\EPair{\E[1]}{\E[2]}}
    {\TPairOrd{\T}{\TU}}
}

% T-Seq
\DefRule {Seq} {
  \HasType
    {\Ctx[1]}
    {\E[1]}
    {\T} \\
  \HasType
    {\Ctx[2]}
    {\E[2]}
    {\TU} \\
  \Unr \T
}{
  \HasType
    {\CtxPat 1 2}
    {\ESeq{\E[1]}{\E[2]}}
    {\TU}
}

% T-Let
\DefRule {Let} {
  \HasType
    {\Ctx[1]}
    {\E[1]}
    {\T[1]} \\
  \HasType
    {\CtxVar x {{\T[1]}} \CtxPatOp \Ctx[2]}
    {\E[2]}
    {\T[2]}
}{
  \HasType
    {\CtxPat 1 2}
    {\ELet x{\E[1]}{\E[2]}}
    {\T[2]}
}

% T-LetUnord
\DefRule {LetUnord} {
  \HasType
    {\Ctx[1]}
    {\E[1]}
    {\TPairUnord {\T[1]}{\T[2]}} \\\\
  \HasType
    {(\CtxVar x{\T[1]} \CtxPar \CtxVar y{\T[2]}) \CtxPatOp \Ctx[2]}
    {\E[2]}
    {\TU}
}{
  \HasType
    {\CtxPat 1 2}
    {\ELetPair xy{\E[1]}{\E[2]}}
    {\TU}
}

% T-LetOrd
\DefRule {LetOrd} {
  \HasType
    {\Ctx[1]}
    {\E[1]}
    {\TPairOrd {\T[1]}{\T[2]}} \\\\
  \HasType
    {(\CtxVar x{\T[1]} \CtxSeq \CtxVar y{\T[2]}) \CtxPatOp \Ctx[2]}
    {\E[2]}
    {\TU}
}{
  \HasType
    {\CtxPat 1 2}
    {\ELetPair xy{\E[1]}{\E[2]}}
    {\TU}
}

% T-Inj
\DefRule {Inj} {
  \HasType
    {\Ctx}
    {\E}
    {\T[i]}
}{
  \HasType
    {\Ctx}
    {\EInj{i} \E}
    {\TSum{\T[1]}{\T[2]}}
}

% T-Case
\DefRule {Case} {
  \HasType
    {\Ctx[1]}
    {\E}
    {\TVariant{\ell:\T[\ell]}}
  \\
  \HasType
    {\CtxVar{\EVar[\ell]}{\T[\ell]} \CtxPatOp \Ctx[2]}
    {\E[\ell]}
    {\TU}
}{
  \HasType
    {\CtxPat 1 2}
    {\ECase [\ell\in L] \E \ECaseBranchDef}
    {\TU}
}

% T-Weaken
\DefRule {Weaken} {
  \Ctx[1] \SubCtx \Ctx[2]
  \\
  \HasType
    {\Ctx[1]}
    {\E}
    {\T}
}{
  \HasType
    {\Ctx[2]}
    {\E}
    {\T}
}

% T-Conv
\DefRule {Conv} {
  \Eff[1] \le \Eff[2]
  \\
  \HasType[\Eff[1]]
    {\Ctx}
    {\E}
    {\T}
  \\
  \T \simeq \TU
}{
  \HasType[\Eff[2]]
    {\Ctx}
    {\E}
    {\TU}
}

% \newcommand*\RuleTMatch{
%   \inferrule[T-MatchLeft]{
%     \HasType[\Eff]
%       {\Ctx[1]}
%       {\E}
%       {\SBranch{\ell:\S[\ell]}{\ell\in L}}
%     \\
%     \HasType[\Eff[\ell]]
%       {\Ctx[2]}
%       {\E[\ell]}
%       {\TArr{\Mob}{\Eff[\ell]}{\S[\ell]}{\T}}
%   }{
%     \HasType[\Eff \EffJoin \EffJoin \{ \Eff[\ell] \}_{\ell\in L}]
%       {\Ctx[1] \CtxSeq \Ctx[2]}
%       {\EMatch{\E}{\ell:\E[\ell]}{\ell\in L}}
%       {\T}
%   }
% }
% 
% \newcommand*\RuleTMatchLin{
%   \inferrule[T-MatchLin]{
%     \HasType[\Eff]
%       {\Ctx[1]}
%       {\E}
%       {\SBranch{\ell:\S[\ell]}{\ell\in L}}
%     \\
%     \HasType[\Eff[\ell]]
%       {\Ctx[2]}
%       {\E[\ell]}
%       {\TArr{\Mob}{\Eff[\ell]}{\S[\ell]}{\T}}
%   }{
%     \HasType[\Eff \EffJoin \EffJoin\{ \Eff[\ell]\}_{\ell\in L}]
%       {\Ctx[1] \CtxPar \Ctx[2]}
%       {\EMatch{\E}{\ell:\E[\ell]}{\ell\in L}}
%       {\T}
%   }
% }

%%% Local

%%% Local Variables:
%%% mode: latex
%%% TeX-master: "../main"
%%% End:

% Unrestricted Types

\def\RuleNamePrefix{U-}
\def\RuleCmdPrefix{RuleU}

% U-Base
\DefRule {Base} {} {
  \Unr \B
}

% U-Unit
\DefRule {Unit} {} {
  \Unr \BUnit
}

% U-Prod
\DefRule {Prod} {
  \Unr \T[1] \\
  \Unr \T[2]
} {
  \Unr \TPair{\T[1]}{\T[2]}
}

% U-Sum
\DefRule {Sum} {
  \Unr \T[1] \\
  \Unr \T[2]
} {
  \Unr \TSum{\T[1]}{\T[2]}
}

% U-Arr
\DefRule {Arr} {} {
  \Unr \TArr{\MobMobile}{\DirNone,\Eff}{\T[1]}{\T[2]}
}

% U-Session
\DefRule {Session} {
  \Skips \S
}{
  \Unr \S
}

% U-Skip
\DefRule {Skip} {} {
  \Unr{\SSkip}
}

% U-Seq
\DefRule {Seq} {
  \Unr \S[1] \\
  \Unr \S[2]
} {
  \Unr \SSeq{\S[1]}{\S[2]}
}

% U-Var
\DefRule {Var} {} {
  \Unr{\SVar}
}

% U-Rec
\DefRule {Rec} {
  \Unr{\S}
} {
  \Unr{(\SRec\SVar\S)}
}

% U-CtxEmpty
\DefRule {CtxEmpty} {} {
  \Unr \CtxEmpty
}

% U-CtxVar
\DefRule {CtxVar} {
  \Unr \T
}{
  \Unr {\CtxVar \EVar \T}
}

% U-CtxSeq
\DefRule {CtxSeq} {
  \Unr \Ctx[1] \\
  \Unr \Ctx[2]
}{
  \Unr \CtxSeq{\Ctx[1]}{\Ctx[2]}
}

% U-CtxPar
\DefRule {CtxPar} {
  \Unr \Ctx[1] \\
  \Unr \Ctx[2]
}{
  \Unr \CtxPar{\Ctx[1]}{\Ctx[2]}
}

%%% Local Variables:
%%% mode: latex
%%% TeX-master: "../main"
%%% End:

% Blocked expressions and processes

\def\RuleCmdPrefix{Rule}
\def\RuleNamePrefix{B-}

\newcommand*\Blocked[1]{\operatorname{\mathsf{Blocked}}\,{#1}}
\newcommand*\Ready[1]{\operatorname{\mathsf{Ready}}\,{#1}}
\newcommand*\BC{\operatorname{\mathsf{BC}}}
\newcommand*\AC{\operatorname{\mathsf{AC}}}

% B-ExpBlocked

\DefRule{ExpValueReady} {} {
  \Ready{\PExp{\CUnit}}
}

\DefRule {ExpConstReady} {
  c \notin \{\CNew, \CFork, \CSplit, \CSSplit, \CDrop, \CDiscard\}
} {
  \Ready{\FC{[\EApp\C\V]}}
}

\DefRule {ParReady} {
  \Ready{\Proc[1]} \\
  \Ready{\Proc[2]}
} {
  \Ready{\PPar{\Proc[1]}{\Proc[2]}}
}

\DefRule {NuReady} {
  \Ready{\Proc}
} {
 \Ready{\PNew{{\BindGroup[1]}}{{\BindGroup[2]}}{\Proc}}
}

% B-ParBlocked
\DefRule {ParBlocked} {
  \Blocked{\Proc[1]} \\
  \Blocked{\Proc[2]}
}{
  \Blocked{\PPar{\Proc[1]}{\Proc[2]}}
}

% B-NuBlockedHere
\DefRule {NuBlockedHere} {
  \{\EVar, \EVary\} \cap \BC(\Proc) \ne  \{\EVar, \EVary\} \\
  \Ready{\Proc}
}{
  \Blocked{\PNew{\BSeq{\EVar}{\BindGroup[1]}}{\BSeq{\EVary}{\BindGroup[2]}}{\Proc}}
}

% B-NuBlockedHereAcq
\DefRule {NuBlockedHereAcq} {
  \EVar \in \AC(\Proc)\\ 
  \BindGroup[i] = \EVary \BindGroup \BSep \EVar \BindGroup* \\
  \Ready{\Proc}
}{
  \Blocked{\PNew{{\BindGroup[1]}}{{\BindGroup[2]}}{\Proc}}
}

% B-NuBlockedThere
\DefRule {NuBlockedThere} {
  \Blocked{\Proc}
}{
  \Blocked{\PNew{{\BindGroup[1]}}{{\BindGroup[2]}}{\Proc}}
}

\newcommand*\BlockedChannelsSend{
  \BC(\FCapp{\EApp\CSend{(\EPair\V\EVar)}}) = \{\EVar\}
}

\newcommand*\BlockedChannelsRecv{
  \BC(\FCapp{\EApp\CRecv\EVar}) = \{\EVar\}
}

\newcommand*\BlockedChannelsSelect{
  \BC(\FCapp{\EApp{\CSelect{k}}\EVar}) = \{\EVar\}
}

\newcommand*\BlockedChannelsBranch{
  \BC(\FCapp{\EApp\CBranch\EVar}) = \{\EVar\}
}

\newcommand*\BlockedChannelsTerm{
  \BC(\FCapp{\EApp\CTerm\EVar}) = \{\EVar\}
}

\newcommand*\BlockedChannelsWait{
  \BC(\FCapp{\EApp\CWait\EVar}) = \{\EVar\}
}

\newcommand*\BlockedChannelsAcquire{
  \BC(\FCapp{\EApp\CAcq\EVar}) =  \emptyset
}

\newcommand*\BlockedChannelsDrop{
  \BC(\FCapp{\EApp\CDrop\EVar}) = \emptyset
}

\newcommand*\BlockedChannelsLSplit{
  \BC(\FCapp{\EApp\CSplit\EVar}) = \emptyset
}

\newcommand*\BlockedChannelsRSplit{
  \BC(\FCapp{\EApp\CSSplit\EVar}) = \emptyset
}

\newcommand*\BlockedChannelsFork{
  \BC(\FCapp{\EApp\CFork\V}) = \emptyset
}

\newcommand*\BlockedChannelsNew{
  \BC(\FCapp{\EApp\CNew\E}) = \emptyset
}

\newcommand*\BlockedChannelsDiscard{
  \BC(\FCapp{\EApp\CDiscard\EVar}) = \emptyset
}

\newcommand*\AcqBlockedChannelsSend{
  \AC(\FCapp{\EApp\CSend{(\EPair\V\EVar)}}) = \emptyset
}

\newcommand*\AcqBlockedChannelsRecv{
  \AC(\FCapp{\EApp\CRecv\EVar}) =  \emptyset
}

\newcommand*\AcqBlockedChannelsSelect{
  \AC(\FCapp{\EApp{\CSelect{k}}\EVar}) =  \emptyset
}

\newcommand*\AcqBlockedChannelsBranch{
  \AC(\FCapp{\EApp\CBranch\EVar}) =  \emptyset
}

\newcommand*\AcqBlockedChannelsTerm{
  \AC(\FCapp{\EApp\CTerm\EVar}) = \emptyset
}

\newcommand*\AcqBlockedChannelsWait{
  \AC(\FCapp{\EApp\CWait\EVar}) =  \emptyset
}

\newcommand*\AcqBlockedChannelsAcquire{
  \AC(\FCapp{\EApp\CAcq\EVar}) = \{\EVar\}
}

\newcommand*\AcqBlockedChannelsDrop{
  \AC(\FCapp{\EApp\CDrop\EVar}) = \emptyset
}

\newcommand*\AcqBlockedChannelsLSplit{
  \AC(\FCapp{\EApp\CSplit\EVar}) = \emptyset
}

\newcommand*\AcqBlockedChannelsRSplit{
  \AC(\FCapp{\EApp\CSSplit\EVar}) = \emptyset
}

\newcommand*\AcqBlockedChannelsFork{
  \AC(\FCapp{\EApp\CFork\V}) = \emptyset
}

\newcommand*\AcqBlockedChannelsNew{
  \AC(\FCapp{\EApp\CNew\E}) = \emptyset
}

\newcommand*\AcqBlockedChannelsDiscard{
  \AC(\FCapp{\EApp\CDiscard\EVar}) = \emptyset
}


%%% Local Variables:
%%% mode: latex
%%% TeX-master: "../main"
%%% End:




\newcommand*\Kind{\kappa}
\newcommand*\KVal{\mathsf{Type}}
\newcommand*\KSess{\mathsf{Session}}
\newcommand*\Qual{Q}
\newcommand*\Prop{q}
\newcommand*\TyCtx{\Phi}
\newcommand*\Entails{\mathrel{\vDash}}
\newcommand*\Dualable{\operatorname{\mathsf{dualable}}}
\newcommand*\NonSkip{\operatorname{\mathsf{nonskip}}}
\newcommand*\Wf{\mathrel{\mathsf{wf}}}
\newcommand*\FunVal[1]{\operatorname{\mathsf{funval}}(#1)}
\newcommand*\TForall[4]{\typecolor{\forall(#1:#2).\,#3\Rightarrow#4}}
\newcommand*\TExistsQ[4]{\typecolor{\exists(#1:#2).\,#3\land#4}}
\newcommand*\TyAbs[3]{\termcolor{\Lambda (#1).\,#2 \Rightarrow #3}}
\newcommand*\Pack[3]{\termcolor{\Keyword{pack}\,\langle #1,#2\rangle\,\Keyword{as}\,#3}}
\newcommand*\Open[4]{\termcolor{\Keyword{open}\,\langle #1,#2\rangle=#3\,\Keyword{in}\,#4}}

\title{Qualified System~F Polymorphism for \thecalculus}
\author{Design note}

\begin{document}
\raggedbottom
\begin{abstract}
This note proposes a conservative, explicitly typed polymorphic extension of
the declarative system of \thecalculus.  Universal and existential types are
qualified by the static predicates already observed by the monomorphic type
system.  Type variables and assumptions about them are kept in an
intuitionistic context separate from the ordered expression-variable context.
Type abstraction, type application, packing, and opening add only
expression-level reductions; the process syntax and process reduction rules
are unchanged.
\end{abstract}
\maketitle

\section{Design summary}

The extension follows explicitly typed System~F.  A universal value is
introduced by a term such as
\[
  \TyAbs{\alpha:\KVal}{\Mobile\alpha}{\V}.
\]
The qualification is checked when the abstraction is instantiated; it is not
a run-time argument.  Existential packages use the same qualifications to
record properties of a hidden witness.

Three features of the monomorphic system require care.  First, a type
variable must distinguish ordinary types from session types.  Second, the
rules inspect more properties than mobility: unrestrictedness, boundedness,
eligibility for channel creation, dualability, and nontriviality of a split
suffix are all relevant.  Third, type abstraction is restricted to an
already-formed value.  The form $\TyAbs{\alpha:\Kind}{\Qual}{\V}$ adds no
closure and delays no computation; its resource properties are inherited from
$\V$.  Finally, fixed points remain restricted to functions, but may pass
through existing type abstractions before reaching the function abstraction.
A fixed point is itself a value and unfolds only when supplied with a type or
value argument.  This guarded form admits polymorphic recursion without
admitting fixed points at arbitrary unrestricted types.

\section{Syntax and static contexts}

Let $\alpha,\beta$ range over polymorphic type variables.  They are distinct
from the recursively bound session variables $X$ of $\mu X.S$.  The extension
adds the following grammar.  Elided cases are inherited unchanged from the
paper, except that the recursive expression $\mu f.\E$ is refined to the
function-shaped value form shown below.

\begin{figure}[H]
  \footnotesize
  \begin{align*}
    \Kind &\grmeq \KVal \grmor \KSess
      \tag{Kinds} \\
    \Prop &\grmeq
       \Unr\T
       \grmor \Mobile\T
       \grmor \Bounded\S
       \grmor \SNew\S
       \\
       &\qquad
       \grmor \Dualable\S
       \grmor \NonSkip\S
       \grmor \T \SEq \TU
      \tag{Atomic qualifications} \\
    \Qual &\grmeq \top \grmor \Prop \grmor \Qual\land\Qual
      \tag{Qualifications} \\
    \TyCtx &\grmeq \emptyset
       \grmor \TyCtx,\alpha:\Kind
       \grmor \TyCtx,\Prop
      \tag{Type-and-constraint contexts} \\
    \T &\grmeq \cdots \grmor \alpha
       \grmor \TForall\alpha\Kind\Qual\T
       \\
       &\qquad
       \grmor \TExistsQ\alpha\Kind\Qual\T
      \tag{Types} \\
    \S &\grmeq \cdots \grmor \alpha \grmor \SDual\S
      \tag{Session types} \\
    \U,\V &\grmeq
       \C
       \grmor \EVar
       \grmor \EAbs\EVar\E
       \grmor \EInj i\V
       \grmor \EPair\U\V
       \\
       &\qquad
       \grmor \TyAbs{\alpha:\Kind}{\Qual}{\V}
       \grmor \Pack\T\V{\TExistsQ\alpha\Kind\Qual\T}
       \grmor \EMu\V\quad(\FunVal\V)
      \tag{Values} \\
    \E &\grmeq \cdots
       \grmor \TyAbs{\alpha:\Kind}{\Qual}{\V}
       \grmor \E{[\T]} \\
       &\qquad
       \grmor \Pack\T\E{\TExistsQ\alpha\Kind\Qual\T}
       \\
       &\qquad
       \grmor \Open\alpha\EVar\E\E
       \grmor \EMu\V
      \tag{Expressions}
  \end{align*}
  \caption{Polymorphic syntax}
  \label{fig:poly-syntax}
\end{figure}

Kinds are deliberately small.  A session has kind $\KSess$ and may also be
used as an ordinary value type; universal and existential types have kind
$\KVal$, not $\KSess$.  Thus quantification can hide a session endpoint inside
a package, but it does not turn the package itself into a session protocol.
Quantification is impredicative at kind $\KVal$: a type argument may itself be
universal or existential.

In particular, $\EMu\V$ is a value, subject to the syntactic judgment
$\FunVal\V$ defined in \cref{sec:poly-recursion}.  The body $\V$ must consist
of zero or more type abstractions followed by one function abstraction.

The context $\TyCtx$ is intuitionistic: its entries admit exchange,
contraction, and weakening.  It contains type-variable declarations and
qualification assumptions about them.  In contrast, $\Ctx$ remains the ordered
resource context of expression variables.  The extended expression typing
judgment is
\[
  \TyCtx\,;\,\Ctx \vdash \E : \T \mid \Eff.
\]
No type variable or qualification is represented by an entry in $\Ctx$.

A type abstraction does not create a closure.  Its body must already be a
value, and evaluation never proceeds underneath $\Lambda$.  Consequently,
universal types need neither a closure annotation nor a latent effect.  Their
mobility and unrestrictedness are derived from the type of the wrapped value.

\section{Qualifications and entailment}

We write $\TyCtx\Entails\Prop$ for constraint entailment and extend it
pointwise to conjunctions.  Subject to the correction concerning unrestricted
sessions below, every predicate rule of the monomorphic system is read as a
Horn clause whose premises are entailments.  Thus, for example,
\RuleName{Te-Acq} remains the rule from $\TyCtx\Entails\Bounded\S$ to
$\TyCtx\Entails\Mobile{(\SSeq\SAcq\S)}$ and is not repeated below.
Qualification contexts are kept in solved form: predicate assumptions are
headed by type variables, while predicates on constructed types are
discharged by the structural predicate rules.  The extension adds only the
postulates in part~(a) of \cref{fig:poly-entailment}; part~(b) records useful
admissible rules.

\begin{figure}[H]
  \small
  \textbf{(a) New postulates augmenting the monomorphic predicates.}
  \begin{mathpar}
    \inferrule*[right=Q-Assume]{\Prop\in\TyCtx}{\TyCtx\Entails\Prop}
    \and
    \inferrule*[right=Q-Unr-Var-Mobile]{
      \alpha:\Kind\in\TyCtx \\
      \TyCtx\Entails\Unr\alpha
    }{
      \TyCtx\Entails\Mobile\alpha
    }
    \and
    \inferrule*[right=Q-New-Var-Dual]{
      \alpha:\KSess\in\TyCtx \\
      \TyCtx\Entails\SNew\alpha
    }{
      \TyCtx\Entails\Dualable\alpha
    }
    \\
    \inferrule*[right=Q-Dual-Atom]{
      \S\in\{\SSkip,\SSend\T,\SRecv\T,\STerm,\SWait,\SVar\}
    }{
      \TyCtx\Entails\Dualable\S
    }
    \and
    \inferrule*[right=Q-Dual-Seq]{
      \TyCtx\Entails\Dualable\S[1] \\
      \TyCtx\Entails\Dualable\S[2]
    }{
      \TyCtx\Entails\Dualable{(\SSeq{\S[1]}{\S[2]})}
    }
    \and
    \inferrule*[right=Q-Dual-Rec]{
      \TyCtx\Entails\Dualable\S
    }{
      \TyCtx\Entails\Dualable{(\SRec\SVar\S)}
    }
    \\
    \inferrule*[right=Q-Dual-Choice]{
      \forall\ell\in L.\ \TyCtx\Entails\Dualable\S[\ell]
    }{
      \TyCtx\Entails
        \Dualable{(\SChoice{\ell:\S[\ell]}{\ell\in L})}
    }
    \and
    \inferrule*[right=Q-Dual-Branch]{
      \forall\ell\in L.\ \TyCtx\Entails\Dualable\S[\ell]
    }{
      \TyCtx\Entails
        \Dualable{(\SBranch{\ell:\S[\ell]}{\ell\in L})}
    }
  \end{mathpar}

  \smallskip
  \textbf{(b) Selected derived rules.}
  \begin{mathpar}
    \inferrule*[right=Q-Unr-Mobile]{\TyCtx\Entails\Unr\T}
      {\TyCtx\Entails\Mobile\T}
    \and
    \inferrule*[right=Q-New-Dual]{\TyCtx\Entails\SNew\S}
      {\TyCtx\Entails\Dualable\S}
    \\
    \inferrule*[right=Q-NonSkip-Ground]{
      \Fv(\S)=\emptyset \\
      \neg(\S\SEq\SSkip)
    }{
      \TyCtx\Entails\NonSkip\S
    }
    \and
    \inferrule*[right=Q-NonSkip-Ctx]{
      \TyCtx\Entails\NonSkip\S
    }{
      \TyCtx\Entails
        \NonSkip{(\SSeq{\S[1]}{\SSeq{\S}{\S[2]}})}
    }
    \\
    \inferrule*[right=Q-Cong]{
      \TyCtx\Entails \T\SEq\TU \\
      \TyCtx\Entails \Prop{[\T/\alpha]}
    }{
      \TyCtx\Entails \Prop{[\TU/\alpha]}
    }
  \end{mathpar}
  \caption{Postulated and derived qualification-entailment rules}
  \label{fig:poly-entailment}
\end{figure}

The forms $\Unr\T$ and $\Mobile\T$ are propositions: $\Unr$ and $\Mobile$ are
predicates on types (equivalently, type qualifiers), not type-level functions
or type constructors.

Here $\SVar$ in \RuleName{Q-Dual-Atom} is a recursively bound session
variable, not a polymorphic variable $\alpha$.  There are deliberately no
dualability rules for $\SDrop$ or $\SAcq$.

The rules in part~(b) are admissible as follows.

\paragraph{Unrestricted implies mobile.}
Proceed by induction over the derivation of $\Unr\T$.  Unit and mobile
functions satisfy the corresponding mobility axioms; products and variants
use the induction hypotheses pointwise.  Universal and existential types use
their respective predicate rules.  The only new leaf is a type variable,
discharged by \RuleName{Q-Unr-Var-Mobile}.  Predicate assumptions are first
simplified to such variable-headed leaves.

\paragraph{New implies dualable.}
Induct over the derivation of $\SNew\S$.  The skip, send, receive, sequence,
choice, branch, and recursion cases match the structural dualability
postulates.  A recursively bound variable uses \RuleName{Q-Dual-Atom}; a
polymorphic type variable uses \RuleName{Q-New-Var-Dual}.  Hence every source
protocol admitted by channel restriction has a defined dual.

\paragraph{Non-skip.}
For a closed session, CFST equivalence decides whether
$\S\SEq\SSkip$.  A negative answer establishes
$\TyCtx\Entails\NonSkip\S$ by the semantic definition of $\NonSkip$ below;
it is not a new logical axiom.  For \RuleName{Q-NonSkip-Ctx}, suppose toward a
contradiction that some satisfying ground substitution makes
$\SSeq{\S[1]}{\SSeq{\S}{\S[2]}}\SEq\SSkip$.  Sequential composition has no
cancellation or annihilator: a sequence is equivalent to $\SSkip$ only if
each factor is.  The middle factor would therefore be equivalent to
$\SSkip$, contradicting the premise.

\paragraph{Congruence.}
Predicate interpretations are defined on equivalence classes of types.  If
$\T\SEq\TU$, substituting either representative into a predicate gives the
same proposition.  This proves \RuleName{Q-Cong}; for $\NonSkip$, in
particular, $\T\SEq\TU$ implies that $\T$ is equivalent to $\SSkip$ exactly
when $\TU$ is.

The predicates have the following roles.

\paragraph{Unrestrictedness.}
$\Unr\T$ is required whenever a value may be weakened, contracted, or ignored,
including sequencing and unrestricted recursion.  Linearity remains the
default, so no positive $\mathsf{linear}$ qualification is needed.  Session
endpoints, including endpoints at $\SSkip$, remain linear.

\paragraph{Mobility.}
$\Mobile\T$ is required for communication and for moving captured values
between processes.  It also controls the equations that identify ordered and
unordered context composition.

\paragraph{Boundedness.}
$\Bounded\S$ is needed separately because it entails mobility of
$\SSeq\SAcq\S$.  Replacing it everywhere by a mobility constraint would lose
the compositional rule used for acquired session suffixes.

\paragraph{Channel creation and duality.}
$\SNew\S$ states that $\S$ is a source protocol accepted by $\CNew$ and by
process restriction.  A suspended dual $\SDual\alpha$ is unavoidable in the
polymorphic type of $\CNew$.  Because duality is undefined for $\SDrop$ and
$\SAcq$, formation of $\SDual\S$ requires $\Dualable\S$.  The existing
$\SNew\S$ predicate entails dualability by the admissible rule above, but the
latter is useful independently for generic code over protocols containing
$\STerm$ or $\SWait$.

\paragraph{Non-skip suffixes.}
The split constants require $\neg(\S[2]\SEq\SSkip)$.  Generic splitting
therefore needs the positive qualification $\NonSkip\S$.  Its semantic meaning
is
\[
  \TyCtx\models\NonSkip\S
  \quad\Longleftrightarrow\quad
  \forall\theta.\ 
    \theta\models\TyCtx \Longrightarrow
    \neg(\theta\S\SEq\SSkip),
\]
where $\theta$ ranges over ground, kind-respecting substitutions.  In
particular, $\NonSkip\alpha$ is a qualification on the type variable $\alpha$:
it limits $\alpha$ to ground instantiations not equivalent to $\SSkip$.  It
follows only from an assumption (possibly transported by equality) or from
stronger constraints; failure to prove $\alpha\SEq\SSkip$ is not evidence for
$\NonSkip\alpha$.  At a type application, a qualification
$\NonSkip\alpha$ therefore requires a proof that the actual session argument
is non-skip.  Closed arguments are handled by
\RuleName{Q-NonSkip-Ground}.

\paragraph{Equality.}
An equality qualification records relationships between type expressions
modulo the existing type and CFST equational theories.  It is useful for generic
code whose protocol decomposition is known through type variables.  All predicate
qualifications are invariant under this equality.

For example, consider a wrapper around local splitting (arrow annotations are
shown explicitly):
\begin{align*}
  \mathsf{splitAs}:{}&
    \typecolor{\forall(\alpha,\beta,\gamma:\KSess).}\\[-.3ex]
  &\quad \typecolor{((\alpha\SEq\SSeq{\beta}{\gamma})\land\NonSkip\gamma)
      \Rightarrow}\\[-.3ex]
  &\quad
      \TArr{\MobMobile}{\DirNone,\EffPure}
        \alpha{\TPairOrd\beta\gamma},\\
  \mathsf{splitAs}{[\alpha]}{[\beta]}{[\gamma]}\,c ={}&
    \EApp{\CSplit{[\beta]}{[\gamma]}}c.
\end{align*}
The instantiated split constant expects $c:\SSeq{\beta}{\gamma}$, whereas the
wrapper receives $c:\alpha$.  The equality qualification is used by
\RuleName{T-Conv} to convert the argument before the application; without the
qualification there is no relation between the three protocol type variables
and the body is not typable.  The $\NonSkip\gamma$ qualification discharges the
independent side condition of local splitting.

There is no separate $\Skips\S$ qualification: the declarative rules that need
this information already use $\S\SEq\SSkip$.  Likewise, $\Ord\T$ is not part of
the proposed interface because no declarative typing rule inspects it.  Effects
and directions remain explicit type annotations rather than properties of type
variables.

The former \RuleName{Q-Skip-Unr} is unsound for endpoint values and is not
adopted.  If $x:\SSkip$ were unrestricted, context contraction would type two
occurrences of $\CDiscard\,x$ using one endpoint assumption.  The first process
reduction consumes the restricted binder for $x$, leaving the second occurrence
without an endpoint.  Weakening has the analogous problem: it can erase a
$\SSkip$ value without performing the process-level discard.  The
session clauses of $\Unr$ in the current monomorphic supplement
(\RuleName{U-Session}, \RuleName{U-Skip}, \RuleName{U-Seq},
\RuleName{U-Var}, and \RuleName{U-Rec}) must therefore be removed, or $\Unr$
must be given a different operational interpretation, before combining the
systems.  Mobility of skip-like sessions remains available through the
existing \RuleName{Te-Skips} rule.

For closed arguments, all atomic qualifications reduce to the existing
decisions, including CFST equivalence for equality and $\NonSkip$.  When a type
contains type variables, entailment uses only assumptions, congruence, and the
structural predicate rules.  In particular, failure to refute a qualification
is not a proof of it.

\section{Kinding and context structure}

The judgment $\TyCtx\vdash\T:\Kind$ is standard.  Selected rules are shown
below.  The judgment $\TyCtx\vdash\Qual\ \Wf$ states that all type variables in
the qualification $\Qual$ are bound in $\TyCtx$ and used according to their
kinds.

\begin{figure}[H]
  \begin{mathpar}
    \inferrule*[right=K-Var]{\alpha:\Kind\in\TyCtx}
      {\TyCtx\vdash\alpha:\Kind}
    \and
    \inferrule*[right=K-Session-Type]{\TyCtx\vdash\S:\KSess}
      {\TyCtx\vdash\S:\KVal}
    \\
    \inferrule*[right=K-Dual]{
      \TyCtx\vdash\S:\KSess \\
      \TyCtx\Entails\Dualable\S
    }{
      \TyCtx\vdash\SDual\S:\KSess
    }
    \\
    \inferrule*[right=K-All]{
      \TyCtx,\alpha:\Kind\vdash\Qual\ \Wf \\
      \TyCtx,\alpha:\Kind,\Qual\vdash\T:\KVal
    }{
      \TyCtx\vdash
        \TForall\alpha\Kind\Qual\T :\KVal
    }
    \\
    \inferrule*[right=K-Exists]{
      \TyCtx,\alpha:\Kind\vdash\Qual\ \Wf \\
      \TyCtx,\alpha:\Kind,\Qual\vdash\T:\KVal
    }{
      \TyCtx\vdash
        \TExistsQ\alpha\Kind\Qual\T :\KVal
    }
  \end{mathpar}
  \caption{Selected kinding rules}
  \label{fig:poly-kinding}
\end{figure}

\begin{figure}[H]
  \footnotesize
  \begin{mathpar}
    \inferrule*[right=QF-Top]{}{
      \TyCtx\vdash\top\ \Wf
    }
    \and
    \inferrule*[right=QF-And]{
      \TyCtx\vdash\Qual[1]\ \Wf \\
      \TyCtx\vdash\Qual[2]\ \Wf
    }{
      \TyCtx\vdash\Qual[1]\land\Qual[2]\ \Wf
    }
    \\
    \inferrule*[right=QF-Unr]{
      \TyCtx\vdash\T:\KVal
    }{
      \TyCtx\vdash\Unr\T\ \Wf
    }
    \and
    \inferrule*[right=QF-Mobile]{
      \TyCtx\vdash\T:\KVal
    }{
      \TyCtx\vdash\Mobile\T\ \Wf
    }
    \and
    \inferrule*[right=QF-Bounded]{
      \TyCtx\vdash\S:\KSess
    }{
      \TyCtx\vdash\Bounded\S\ \Wf
    }
    \\
    \inferrule*[right=QF-New]{
      \TyCtx\vdash\S:\KSess
    }{
      \TyCtx\vdash\SNew\S\ \Wf
    }
    \and
    \inferrule*[right=QF-Dualable]{
      \TyCtx\vdash\S:\KSess
    }{
      \TyCtx\vdash\Dualable\S\ \Wf
    }
    \and
    \inferrule*[right=QF-NonSkip]{
      \TyCtx\vdash\S:\KSess
    }{
      \TyCtx\vdash\NonSkip\S\ \Wf
    }
    \\
    \inferrule*[right=QF-Eq]{
      \TyCtx\vdash\T:\Kind \\
      \TyCtx\vdash\TU:\Kind
    }{
      \TyCtx\vdash\T\SEq\TU\ \Wf
    }
  \end{mathpar}
  \caption{Qualification well-formedness ($\TyCtx\vdash\Qual\ \Wf$)}
  \label{fig:qualification-well-formedness}
\end{figure}

Thus $\Unr$ and $\Mobile$ apply to value types; $\Bounded$, $\SNew$,
$\Dualable$, and $\NonSkip$ apply to session types; and equality relates two
types of a common kind.  By \RuleName{K-Var}, the kinding premises are
derivable only when every type variable occurring in the qualification has a
declaration in $\TyCtx$.

All inherited type constructors receive their expected kinding rules.  In
particular, sequential composition and recursive session formation require
$\KSess$ operands, while the payload of send or receive only requires
$\KVal$.  The qualifier entries of $\TyCtx$ are propositions, not kinds.

Context equality and the subcontext preorder must be indexed by $\TyCtx$:
\[
  \TyCtx\vdash\Ctx[1]=\Ctx[2]
  \qquad\text{and}\qquad
  \TyCtx\vdash\Ctx[1]\SubCtx\Ctx[2].
\]
The equations from the paper are unchanged, but their side conditions use
entailment.  For example, $\Unr\alpha\in\TyCtx$ permits an assumption
$x:\alpha$ to contract, and $\Mobile\alpha\in\TyCtx$ permits it to move between
ordered and unordered context composition.  This parameterization is
essential: keeping only $\Ctx$ in the judgment would make qualifications on
type variables invisible to the ordered context algebra.

\section{Universal types}

\begin{figure}[H]
  \begin{mathpar}
    \inferrule*[right=T-TAbs]{
      \TyCtx,\alpha:\Kind,\Qual\,;\,\Ctx
        \vdash \V:\T\mid\EffPure \\
      \TyCtx,\alpha:\Kind\vdash\Qual\ \Wf \\
      \alpha\notin\Fv(\Ctx)
    }{
      \TyCtx\,;\,\Ctx \vdash
        \TyAbs{\alpha:\Kind}{\Qual}{\V}
        :\TForall\alpha\Kind\Qual\T
        \mid\EffPure
    }
    \\
    \inferrule*[right=T-TApp]{
      \TyCtx\,;\,\Ctx\vdash \E:
        \TForall\alpha\Kind\Qual\T
        \mid\Eff \\
      \TyCtx\vdash\TU:\Kind \\
      \TyCtx\Entails\Qual{[\TU/\alpha]}
    }{
      \TyCtx\,;\,\Ctx\vdash
        \E{[\TU]}:\T{[\TU/\alpha]}\mid\Eff
    }
  \end{mathpar}
  \caption{Universal introduction and elimination}
  \label{fig:poly-universal-typing}
\end{figure}

The abstraction rule applies only when the body is already a value.  It checks
that value under a fresh type variable and the declared qualification.  Type
application introduces no latent computation: it evaluates its operator and
then substitutes the type argument into the wrapped value.  Both rules are
therefore pure apart from the effect needed to evaluate the operator of the
type application.

The new predicate rules for universal types are
\begin{mathpar}
  \inferrule{
    \TyCtx,\alpha:\Kind,\Qual\Entails\Mobile\T
  }{
    \TyCtx\Entails
      \Mobile{(\TForall\alpha\Kind\Qual\T)}}
  \and
  \inferrule{
    \TyCtx,\alpha:\Kind,\Qual\Entails\Unr\T
  }{
    \TyCtx\Entails
      \Unr{(\TForall\alpha\Kind\Qual\T)}}.
\end{mathpar}
These rules mirror the existential predicate rules below.  Their
soundness relies on value inversion: a value of mobile, respectively
unrestricted, type can use only a mobile, respectively unrestricted,
expression context.  This argument would fail for an arbitrary expression,
which is why \RuleName{T-TAbs} requires the body to be a value.

\section{Existential types}

Existential introduction supplies an explicit witness.  Elimination exposes a
fresh type variable, its qualification, and the packed value.  The
result type and the surrounding expression-variable context may not mention
the hidden variable.

\begin{figure}[H]
  \begin{mathpar}
    \inferrule*[right=T-Pack]{
      \TyCtx\vdash\TU:\Kind \\
      \TyCtx\Entails\Qual{[\TU/\alpha]} \\
      \TyCtx\,;\,\Ctx\vdash
        \E:\T{[\TU/\alpha]}\mid\Eff
    }{
      \TyCtx\,;\,\Ctx\vdash
      \Pack\TU\E{\TExistsQ\alpha\Kind\Qual\T}
      :\TExistsQ\alpha\Kind\Qual\T\mid\Eff
    }
    \\
    \inferrule*[right=T-Open]{
      \TyCtx\,;\,\Ctx[1]\vdash
        \E[1]:\TExistsQ\alpha\Kind\Qual\T\mid\Eff \\
      \TyCtx,\alpha:\Kind,\Qual\,;\,
        (\CtxVar\EVar\T\CtxPatOp\Ctx[2])
        \vdash\E[2]:\TU\mid\Eff \\
      \alpha\notin\Fv(\Ctx[2],\TU)
    }{
      \TyCtx\,;\,(\Ctx[1]\CtxPatOp\Ctx[2])\vdash
        \Open\alpha\EVar{\E[1]}{\E[2]}:\TU\mid\Eff
    }
  \end{mathpar}
  \caption{Existential introduction and elimination}
  \label{fig:poly-existential-typing}
\end{figure}

As in the paper's let rule, $\CtxPatOp$ ranges over ordered and unordered
context composition and is chosen consistently in the premise and conclusion.
The qualification is available only while checking the body of the open.

An existential package has the run-time resources of its payload.  Its
mobility and unrestrictedness are therefore derived under the bound witness
type variable and its qualification:
\begin{mathpar}
  \inferrule{
    \TyCtx,\alpha:\Kind,\Qual\Entails\Mobile\T
  }{
    \TyCtx\Entails\Mobile{(\TExistsQ\alpha\Kind\Qual\T)}
  }
  \and
  \inferrule{
    \TyCtx,\alpha:\Kind,\Qual\Entails\Unr\T
  }{
    \TyCtx\Entails\Unr{(\TExistsQ\alpha\Kind\Qual\T)}.
  }
\end{mathpar}
For example,
\[
  \TExistsQ\alpha\KVal{\Mobile\alpha}\alpha
\]
is mobile, while $\TExistsQ\alpha\KVal\top\alpha$ is not.  Similarly,
$\TExistsQ\alpha\KSess{\Bounded\alpha}{(\SSeq\SAcq\alpha)}$ is mobile by
the acquired-session rule.

\section{Lifting the monomorphic rules}

Every declarative judgment in the paper is prefixed by $\TyCtx$.  The shape of
the inherited rules is otherwise unchanged.  Predicate and equality side
conditions become entailments.  Representative replacements are
\[
\begin{array}{lll}
  \RuleName{T-Seq}:& \TyCtx\Entails\Unr\T,
    &\text{the discarded result is unrestricted};\\
  \RuleName{CT-Send}/\RuleName{CT-Recv}:&
    \TyCtx\Entails\Mobile\T,
    &\text{the payload crosses a process boundary};\\
  \RuleName{CT-LSplit}/\RuleName{CT-RSplit}:&
    \TyCtx\Entails\NonSkip\S[2],
    &\text{the residual suffix is nontrivial};\\
  \RuleName{TP-Res}:& \TyCtx\Entails\SNew\S,
    &\text{restriction uses a source protocol};\\
  \RuleName{T-Conv}:& \TyCtx\Entails\T\SEq\TU,
    &\text{conversion may use abstract equations}.
\end{array}
\]
The abstraction rules for ordinary functions continue to test mobility or
unrestrictedness of their captured expression context through entailment.
Process typing becomes $\TyCtx\,;\,\Ctx\vdash\Proc$; binder-group typing is
similarly parameterized.  Type variables never enter binder groups at run
time.

The type-indexed primitive families can now be presented as explicitly
polymorphic constants.  Representative signatures are
\begin{align*}
  \CSend :{}&
    \TForall\alpha\KVal{\Mobile\alpha}
      {\TArr{\MobMobile}{\DirNone,\EffImpure}
        {\TPairUnord\alpha{\SSend\alpha}}{\BUnit}},\\
  \CRecv :{}&
    \TForall\alpha\KVal{\Mobile\alpha}
      {\TArr{\MobMobile}{\DirNone,\EffImpure}{\SRecv\alpha}{\alpha}},\\
  \CNew :{}&
    \TForall\alpha\KSess{\SNew\alpha}
      {\TArr{\MobMobile}{\DirNone,\EffPure}{\BUnit}
      {\TPairUnord
        {\SSeq\SAcq{\SSeq\alpha\STerm}}
        {\SSeq\SAcq{\SSeq{\SDual\alpha}\SWait}}}},\\
  \CSplit :{}&
    \TForall\alpha\KSess\top
      {\TForall\beta\KSess{\NonSkip\beta}
        {\TArr{\MobMobile}{\DirUnord,\EffPure}
          {\SSeq\alpha\beta}{\TPairOrd\alpha\beta}}},\\
  \CSSplit :{}&
    \TForall\alpha\KSess\top
      {\TForall\beta\KSess{\NonSkip\beta}
        {\TArr{\MobMobile}{\DirUnord,\EffPure}
          {\SSeq\alpha\beta}
          {\TPairUnord{\SSeq\alpha\SDrop}{\SSeq\SAcq\beta}}}}.
\end{align*}
Thus $\CSend_{\T}$ is written $\CSend{[\T]}$, and similarly for the other
formerly indexed constants.  The existing monomorphic constant rules are the
instances of these signatures.  Choice and branch constants can quantify the
finite family $(\S[\ell]:\KSess)_{\ell\in L}$ componentwise; no row
polymorphism is required for the core extension. \os{What does this sentence mean?}

\section{Recursion}
\label{sec:poly-recursion}

Recursion is factored from unrestricted function introduction.  This
factorization admits polymorphic recursion while retaining the operational
restriction that every fixed point is guarded by a function abstraction.
First, writing
\[
  A = \TArr{\MobMobile}{\DirNone,\Eff}\T\TU,
\]
an unrestricted function is introduced independently of recursion:
\begin{mathpar}
  \inferrule*[right=T-AbsUnr]{
    \TyCtx\,;\,
      (\CtxVar{\EVar}{\T} \CtxPar \Ctx)
      \vdash\E:\TU\mid\Eff \\
    \TyCtx\Entails\Unr\Ctx
  }{
    \TyCtx\,;\,\Ctx\vdash
      \EAbs\EVar\E:A\mid\EffPure
  }
\end{mathpar}
This is the abstraction component of the paper's existing
\RuleName{T-AbsUnr}; the recursive binder is moved to a separate rule.

The following syntactic judgment recognizes values whose erasure is a
function.  Type abstractions may precede the lambda, but no pair, package, or
other value constructor may intervene.
\begin{mathpar}
  \inferrule*[right=FV-Abs]{}{
    \FunVal{\EAbs\EVar\E}
  }
  \and
  \inferrule*[right=FV-TAbs]{
    \FunVal\V
  }{
    \FunVal{\TyAbs{\alpha:\Kind}{\Qual}{\V}}
  }
\end{mathpar}
The fixed-point rule ranges over exactly these values:
\begin{mathpar}
  \inferrule*[right=T-Fix]{
    \TyCtx\,;\,(\CtxVar f\T \CtxPar \Ctx)
      \vdash\V:\T\mid\EffPure \\
    \TyCtx\Entails\Unr\T \\
    \FunVal\V
  }{
    \TyCtx\,;\,\Ctx\vdash
      \EMu\V:\T\mid\EffPure
  }.
\end{mathpar}
The conclusion $\EMu\V$ is a value by the grammar in
\cref{fig:poly-syntax}; \RuleName{T-Fix} does not trigger an unfolding step.
Because the final lambda is checked by \RuleName{T-AbsUnr}, its captured
context, including $f:\T$, must be unrestricted.  The explicit
$\Unr\T$ premise records the corresponding requirement on the recursive
variable.  The original rule is recovered by taking
$\V=\EAbs\EVar\E$.

For polymorphic recursion, take
\[
  \V =
    \TyAbs{\alpha:\Kind}{\Qual}{\EAbs\EVar\E}
  \qquad\text{and}\qquad
  \T = \TForall\alpha\Kind\Qual A.
\]
The premise of \RuleName{T-Fix} checks $\V$ with $f:\T$ in scope.  Hence the
body $\E$ may invoke $f$ at any admissible type argument, including one
different from $\alpha$.  Further type abstractions are handled by repeated
applications of \RuleName{FV-TAbs}; no separate big-lambda form is required.

The declarative rule may leave $\T$ implicit.  An algorithmic presentation
should instead use an annotation such as $\EMu[(f:\T)]\V$, or require an
equivalent surrounding ascription, because the complete polymorphic type of a
recursive variable must be supplied rather than inferred circularly.  The
annotation supplies only the self type; its type variables remain bound by the
ordinary $\Lambda$ forms in $\V$.

\section{Expression dynamics}

The complete value grammar is given in \cref{fig:poly-syntax}.  In particular,
a fixed point satisfying $\FunVal\V$ is a value and does not reduce in
isolation.  The evaluation contexts evaluate the operator of a type
application, the payload of a package, and the scrutinee of an open:
\begin{align*}
  \EC &\grmeq \cdots
    \grmor \EC{[\T]}
    \grmor \Pack\T\EC{\TExistsQ\alpha\Kind\Qual\T}
    \grmor \Open\alpha\EVar\EC\E.
\end{align*}
The new axiomatic reductions are
\begin{mathpar}
  \inferrule*[right=E-TApp]{}{
    (\TyAbs{\alpha:\Kind}{\Qual}{\V}){[\T]}
      \EReduce \V{[\T/\alpha]}
  }
  \and
  \inferrule*[right=E-Open]{}{
    \Open\alpha\EVar
      {\Pack\T\V{\TExistsQ\beta\Kind\Qual\TU}}
      \E
    \EReduce
      \E{[\T/\alpha]}{[\V/\EVar]}
  }.
\end{mathpar}
Qualification annotations have no dynamic evidence.  Type abstraction adds no
closure and suspends no computation: its body is already a value, and
\RuleName{E-TApp} only performs type substitution in that value.  No evaluation
context descends underneath $\Lambda$.  Type arguments, qualifications, and
the $\Lambda$ wrapper can therefore be erased without introducing a run-time
thunk.

Unfolding is demand driven.  Either a type argument or a value argument
triggers one unfolding step, while the argument remains in place:
\begin{mathpar}
  \inferrule*[right=E-FixTApp]{}{
    \bigl(\EMu\V\bigr){[\T]}
      \EReduce
    \bigl(\Subs\V{\EMu\V}f\bigr){[\T]}
  }
  \and
  \inferrule*[right=E-FixApp]{}{
    \EApp{(\EMu\V)}\U
      \EReduce
    \EApp{(\Subs\V{\EMu\V}f)}\U
  }
\end{mathpar}
The rules do not inspect the shape of $\V$ beyond the restriction imposed by
the value grammar.  Typing inversion guarantees that a fixed point receiving a
type argument unfolds to a type abstraction, whereas one receiving a value
argument unfolds to a lambda.  The ordinary \RuleName{E-TApp} or
\RuleName{E-App} rule can therefore take the next step.  Substitution retains
the original fixed point at recursive occurrences, so they keep the complete
polymorphic type and may instantiate it differently.  No reduction applies to
a fixed point that is merely stored, passed, or returned.

No process form and no process reduction axiom is added.  The existing
\RuleName{R-Exp} rule already lifts every expression reduction through a
process-local evaluation context.  Communication, splitting, restriction,
drop/acquire synchronization, process congruence, and all structural closure
rules are literally unchanged.  This is the sense in which polymorphism is
operationally orthogonal to the process semantics.

\section{Metatheoretic obligations}

The extension adds the standard lemmas for qualified System~F.

\begin{enumerate}
\item \emph{Entailment substitution.}  If
  $\TyCtx,\alpha:\Kind,\Qual\Entails\Prop$,
  $\TyCtx\vdash\TU:\Kind$, and
  $\TyCtx\Entails\Qual{[\TU/\alpha]}$, then
  $\TyCtx\Entails\Prop{[\TU/\alpha]}$.
\item \emph{Type substitution.}  Typing, kinding, context equality, and the
  subcontext relation are preserved by a type substitution satisfying the
  qualification of its binder.
\item \emph{Existential escape.}  The side condition on \RuleName{T-Open}
  ensures that opening a package cannot expose its witness in the result type
  or in an outer resource assumption.
\item \emph{Fixed-point canonical forms.}  The $\FunVal\V$ judgment is
  preserved by type and term substitution.  Typing inversion shows that the
  body of a fixed-point value at universal type starts with a type
  abstraction, whereas the body of one at arrow type is a lambda.  Hence a
  well-typed fixed point reduces precisely when the corresponding type or
  value argument is present.
\item \emph{Preservation and progress.}  Type application and package opening
  follow from type substitution and ordinary term substitution.  The
  fixed-point application cases additionally use the substitution lemma for
  the unrestricted recursive variable.  Universal canonical forms are type
  abstractions or fixed-point values whose bodies start with a type
  abstraction; arrow canonical forms are lambdas or fixed-point values whose
  bodies are lambdas.  Existential canonical forms are packages.
\item \emph{Process safety.}  Process preservation and progress require no
  new process cases.  Their expression case invokes the extended expression
  theorems.
\end{enumerate}

For an algorithmic system, $\TyCtx$ should remain separate from the paper's
generated unification constraints $\CSet$.  Type variables declared in
$\TyCtx$ are parameters and are never solved; the variables generated in
$\CSet$ are unification variables and may be solved.  Explicit abstraction and
application make quantifier placement and instantiation inputs to checking,
while qualification entailment reuses the existing predicate checks and CFST
equivalence procedure.  Equality assumptions over type variables are used by
congruence and conversion, not solved by unification.  The algorithmic
fixed-point rule additionally checks the supplied type of the recursive
variable; it does not infer a polymorphic recursive type.

\section{Recommended core}

The smallest useful version contains both kinds, universal and existential
types, the qualifications
\[
  \Unr\T,\quad \Mobile\T,\quad \Bounded\S,\quad \SNew\S,\quad
  \Dualable\S,\quad \NonSkip\S,\quad \T\SEq\TU,
\]
together with the value restriction on type abstraction and the
function-shape restriction on fixed points.  Of these,
$\Dualable$ can be hidden from surface syntax if suspended duals are generated
only by the polymorphic type of $\CNew$; it is still needed internally to make
$\SDual\alpha$ well formed.  Equality qualifications may likewise be omitted
from an initial implementation if protocol refinements involving type variables
are not required.  The remaining five predicates are directly observed by the
declarative rules and should not be collapsed into mobility alone.

\end{document}

%%% Local Variables:
%%% mode: latex
%%% TeX-master: t
%%% End:
